% ===================================================================
% FST-II: Chemische Stabilit\"at und autokatalytische Selektion
% Zieljournal: Origins of Life and Evolution of Biospheres / J. Theor. Biol.
% Version: 2.0 (integriertes Quellmaterial)
% ===================================================================

\documentclass[12pt,a4paper]{article}

% Packages
\usepackage[utf8]{inputenc}
\usepackage[T1]{fontenc}
\usepackage[ngerman]{babel}
\usepackage[protrusion=true,expansion=false]{microtype}
\usepackage{amsmath,amssymb,amsthm}
\usepackage{physics}
\usepackage{graphicx}
\usepackage[unicode=true,pdfencoding=auto,psdextra]{hyperref}
\usepackage{booktabs}
\usepackage[margin=2.5cm]{geometry}
\usepackage{natbib}
\usepackage{cleveref}
\usepackage{enumitem}
\usepackage{tabularx}

\setlength{\emergencystretch}{3em}
\makeatletter
% Keep EN and DE hyperlink destinations distinct after PDF merge.
\renewcommand*{\HyperDestNameFilter}[1]{fst2chem-de-#1}
\makeatother

% Theorem environments
\newtheorem{proposition}{Proposition}
\newtheorem{conjecture}{Conjecture}
\newtheorem{definition}{Definition}
\newtheorem{theorem}{Theorem}
\newtheorem{lemma}{Lemma}
\theoremstyle{remark}
\newtheorem{remark}{Remark}

% Title
\title{Funktionale Stabilit\"atstheorie II: Chemische Stabilit\"at und autokatalytische Selektion}

\author{Lukas Geiger\\
\textit{Unabh\"angiger Forscher, Bernau, Deutschland}\\
ORCID: \url{https://orcid.org/0009-0005-7296-1534}}

\date{Mai 2026}

\hypersetup{
    pdftitle={Funktionale Stabilitätstheorie II: Chemische Stabilität und autokatalytische Selektion},
    pdfauthor={Lukas Geiger},
    pdfsubject={Präbiotische Chemie, Autokatalyse und Funktionale Stabilitätstheorie},
    pdfkeywords={Autokatalyse, Hyperzyklus, Spieltheorie, Nash-Gleichgewicht, Chiralität, Ursprung des Lebens, Entropieproduktion, Dynamic Kinetic Stability, Replikator-Dynamik}
}

\begin{document}

\maketitle

\begin{abstract}
Dieser Artikel entwickelt einen spieltheoretischen Rahmen f\"ur die pr\"abiotische Chemie und untersucht, wie autokatalytische Netzwerke, Hyperzyklen und Chiralit\"atsselektion als Nash-Gleichgewichts-artige Strukturen thermodynamischer Spiele interpretiert werden k\"onnen. Aufbauend auf Dynamic Kinetic Stability (DKS) und der England-Ungleichung aus der statistischen Mechanik des Nichtgleichgewichts schlagen wir vor, das Maximum Entropy Production Principle (MEPP) als Stabilit\"atsfilter zu lesen, der die lebensf\"ahige Mannigfaltigkeit replizierender Systeme abgrenzt, statt als kausalen Treiber molekularer Evolution. Empirische Motivation liefern Azoarcus-Ribozym-Experimente, in denen sich RNA-Replikationsdynamiken auf klassische spieltheoretische Payoff-Matrizen abbilden lassen, darunter Kooperation, Dominanz und Gefangenendilemma. Der Rahmen modelliert Homochiralit\"at als spontane Symmetriebrechung in einem entarteten Nash-Gleichgewicht, konsistent mit Viedma-Ripening-Experimenten. R\"aumliche Erweiterungen \"uber Reaktions-Diffusions-Systeme legen nahe, dass Spiralwellen als verteilte evolution\"ar stabile Strategie (Distributed Evolutionarily Stable Strategy, DESS) gedeutet werden k\"onnen, die Hyperzyklen gegen parasit\"are Invasion sch\"utzt. Eine Resolventenstabilit\"ats-Perspektive fasst strukturelle Selektion zweiter Ordnung als Kandidatenmechanismus, der persistente von transienten Konfigurationen unterscheidet. Es werden sechs testbare Vorhersagen formuliert. Die Grenzen des Rahmens---insbesondere der umstrittene Status von MEPP fern vom Gleichgewicht und die Beschr\"ankung auf ein einziges experimentelles System---werden ausdr\"ucklich benannt.
\end{abstract}

\textbf{Schl\"usselw\"orter:} Autokatalyse, Hyperzyklus, Spieltheorie, Nash-Gleichgewicht, Chiralit\"at, Ursprung des Lebens, Entropieproduktion, Dynamic Kinetic Stability, Replikator-Dynamik

\clearpage

% ===================================================================
\section{Einleitung}
\label{sec:introduction}
% ===================================================================

Der Ursprung des Lebens repr\"asentiert eine der tiefgreifendsten Herausforderungen der modernen Wissenschaft. Traditionelle Ans\"atze zu diesem Problem waren weitgehend deskriptiv und konzentrierten sich auf die Katalogisierung molekularer Bausteine und die Analyse isolierter linearer Reaktionspfade in mutma\ss{}lichen pr\"abiotischen Umgebungen \citep{MillerUrey}. W\"ahrend diese klassischen Ans\"atze grundlegende Evidenz f\"ur die abiotische Synthese organischer Molek\"ule lieferten, konnten sie den entscheidenden qualitativen Sprung nicht erkl\"aren: die spontane Entstehung systemischer Komplexit\"at, Selbstreplikation und evolution\"arer Dynamik aus einer stochastischen Mischung unbelebter Materie.

Der FST-II-Rahmen (Functional Stability Theory II---``Fractal'' in dem Sinne, dass dieselbe Stabilit\"atslogik zweiter Ordnung \"uber physikalische Skalen hinweg wiederkehrt) schl\"agt eine komplement\"are Perspektive vor. Er erweitert den spieltheoretischen Formalismus, der in FST-I \citep{FSTI} eingef\"uhrt wurde, von der Teilchenphysik auf chemische Skalen, aufbauend auf den thermodynamischen Grundlagen von Englands statistischer Physik der Selbstreplikation \citep{England2013}. Er verl\"asst die reduktionistische Ebene reiner molekularer Katalogisierung und modelliert den \"Ubergang von pr\"abiotischer Chemie zu biologischen Systemen als einen m\"oglichen Stabilit\"atsselektionsprozess: Chemische Konfigurationen werden durch Nicht-Gleichgewichtsthermodynamik eingeschr\"ankt, w\"ahrend ihre Persistenz mit dem Vokabular der evolution\"aren Spieltheorie analysiert wird.

Im Kernpostulat von FST-II steht die Pr\"amisse, dass pr\"abiotische Ph\"anomene h\"oherer Ordnung---wie die Etablierung der Autokatalyse, die Entstehung informationstragender Hyperzyklen \citep{Eigen1979} und die globale Symmetriebrechung der biologischen Homochiralit\"at---nicht ausschlie\ss{}lich als zuf\"allige historische Unf\"alle oder statistische Ausrei\ss{}er behandelt werden m\"ussen. Vielmehr k\"onnen diese Ph\"anomene statistisch robuste makroskopische Attraktoren in einem hochdimensionalen chemischen Phasenraum repr\"asentieren \citep{ReplicatorEntropy}. Konkret formuliert die Theorie diese Zust\"ande als Nash-Gleichgewichts-artige Strukturen in einem thermodynamischen Mehrspieler-Spiel, in dem einzelne Molek\"ule, pr\"abiotische Oligomere und ganze Reaktionsnetzwerke \emph{so modelliert werden, als ob} sie ``Spieler'' w\"aren, deren ``Strategien'' durch molekulare Erkennungsmechanismen, sterische Konformationen und katalytische Affinit\"aten definiert sind \citep{CGT2018,CGTMacmillan}.

% ===================================================================
\section{Strukturelle Ausrichtung: FST-II als Pattern-A-Instanz}
\label{sec:structural-alignment}
% ===================================================================

Der hier entwickelte chemische Stabilit\"atsrahmen (DKS) wird als vorgeschlagene Instanz der \textbf{Pattern-A-Stabilit\"atslogik} behandelt, die das FST-\"Okosystem vereinigt (siehe \cite{FST-Unified2026}). Pattern A ist eine wiederkehrende mathematische Struktur \"uber FST hinweg: ein System zeigt Entartung in f\"uhrender Ordnung (viele Konfigurationen sind in erster N\"aherung gleicherma\ss{}en lebensf\"ahig), und die Selektion erfolgt in zweiter Ordnung durch Kopplung zwischen entarteten Unterr\"aumen. Die j\"ungste \textbf{eichtheoretische Reformulierung der Riemannschen Vermutung} (siehe \cite{GeigerRH2026c}) liefert eine strukturelle Vorlage f\"ur diese Architektur; dieser Abschnitt bildet diese Vorlage auf die pr\"abiotische Chemie ab. Leser, die prim\"ar an der chemischen Spieltheorie interessiert sind, k\"onnen ohne Verlust der Kontinuit\"at zu Abschnitt~\ref{sec:thermodynamics} vorspringen.

In diesem Kontext werden pr\"abiotische autokatalytische Netzwerke wie folgt verstanden:
\begin{enumerate}[label=\alph*)]
    \item \textbf{Analytische Rang-Endlichkeit (Null-Tail):} Der riesige kombinatorische Raum m\"oglicher Molek\"ulsequenzen bleibt handhabbar, weil nur eine endliche Teilmenge von ``Knowlecules'' replikative Stabilit\"at erreicht. Der unorganisierte chemische Rest wird durch die entropische Einschr\"ankung dissipiert, was verhindert, dass das System in einen ungerichteten Zufallsprozess zerf\"allt.
    \item \textbf{Endliche Negativit\"at als Eichsignal:} Parasit\"are Sequenzen und lokale Instabilit\"aten sind das chemische Analogon der endlichen Negativit\"at im arithmetischen Kern. Sie sind keine Defekte, sondern Signale einer fehlenden organisatorischen Eichung---die durch Kompartimentierung und r\"aumliche Symmetriebrechung (Spiralwellen) aufgel\"ost wird.
    \item \textbf{Eingeschr\"ankte funktionale Positivit\"at (Chemische ESS):} Die ESS eines chemischen Netzwerks ist der ``eichstabile'' Zustand, in dem das totale Dissipationsfunktional resolventenstabilisiert ist (nicht notwendig global maximiert; siehe die Stabilit\"atsfilter-Interpretation in Abschnitt~\ref{sec:mepp}). So wie die RH auf eine Rang-2-Stabilisierung reduziert wird, sind chemische Homochiralit\"at und Hyperzyklen das Ergebnis einer endlichen Menge von Einschr\"ankungen, die die globale Positivit\"at der Replikationsdynamik gew\"ahrleisten.
\end{enumerate}

Diese Ausrichtung erhebt FST-II von einer inspirierten Analogie zu einer strukturellen Parallele von ``Funktionale Positivit\"at unter Einschr\"ankung''.

\paragraph{Terminologie.} Im gesamten Artikel bezeichnet \emph{resolvent-stabil} einen Fixpunkt, bei dem die linearisierte Dynamik strikt negative Realteil-Eigenwerte in Nicht-Eich-Richtungen aufweist---d.\,h.\ asymptotische Stabilit\"at im ODE-Sinne (Abschnitt~\ref{sec:mepp}). Im chemischen Kontext ist dies \"aquivalent zur Standard-asymptotischen Stabilit\"at, wie sie in der chemischen Kinetik seit Ljapunow (1892) definiert ist. Wir \"ubernehmen das ``Resolvent''-Vokabular als \emph{terminologische Br\"ucke} zum FST-RH-Rahmen, wo es zus\"atzlichen spektraltheoretischen Gehalt tr\"agt; im vorliegenden chemischen Rahmen f\"ugt es keinen neuen mathematischen Gehalt \"uber asymptotische Stabilit\"at hinaus, hebt aber die strukturelle Parallele \"uber Skalen hinweg hervor. Man beachte, dass im FST-RH-Rahmen \citep{GeigerRH2026b} ``Positivit\"at'' sich auf die Hesse-Matrix eines Variationsfunktionals bezieht (positive Eigenwerte $=$ Maximum), w\"ahrend sie sich hier auf die Jacobi-Matrix der Ratengleichungen bezieht (negative Realteile $=$ R\"uckkehr in den station\"aren Zustand). Beide Konventionen kodieren strukturelle Stabilit\"at; der Vorzeichenunterschied reflektiert den \"Ubergang von einer variationellen zu einer dynamischen Formulierung.

% ===================================================================
\section{Thermodynamische Grundlagen: Dynamic Kinetic Stability}
\label{sec:thermodynamics}
% ===================================================================

Um chemische und pr\"abiotische Reaktionsnetzwerke spieltheoretisch zu modellieren, muss zun\"achst die thermodynamische Makroumgebung pr\"azise definiert werden.

\subsection{Die Dichotomie der Stabilit\"atskonzepte}

Das klassische Verst\"andnis chemischer Reaktionen in geschlossenen Systemen basiert auf dem Antrieb zu einem Zustand maximaler Systementropie und minimaler freier Energie---der Bedingung thermodynamischer Stabilit\"at \citep{Pross2012}. Jedoch existieren Leben und seine pr\"abiotischen Vorl\"auferstrukturen niemals in isolierten, geschlossenen Systemen. Sie sind obligatorisch offene Systeme, die weit vom thermodynamischen Gleichgewicht operieren und auf permanenten Energie- und Materiefluss angewiesen sind \citep{Bioenergetics}.

Die FST-II-Theorie beruht auf dem Konzept der Dynamic Kinetic Stability (DKS) \citep{Pross2012}:

\begin{definition}[Dynamic Kinetic Stability]
DKS beschreibt das Verhalten hochenergetischer Systeme, die weit vom thermodynamischen Gleichgewicht entfernt lokalisiert sind und ihre strukturelle Integrit\"at, innere Ordnung und funktionale Persistenz \"uber die Zeit aufrechterhalten---nicht durch energetische Stasis oder Isolation, sondern paradoxerweise durch den kontinuierlichen Durchsatz und Umsatz von Materie und Energie aus ihrer Umgebung.
\end{definition}

\begin{center}
\begin{tabularx}{\textwidth}{>{\raggedright\arraybackslash}p{0.20\textwidth}>{\raggedright\arraybackslash}X>{\raggedright\arraybackslash}X}
\toprule
\textbf{Dimension} & \textbf{Thermodynamische Stabilit\"at} & \textbf{Dynamic Kinetic Stability} \\
\midrule
Systemzustand & Statisches Gleichgewicht & Nicht-Gleichgewichts-Steady-State (NESS) \\
Energieniveau & Absolutes Minimum & Hochenergetisch (durch Fluss aufrechterhalten) \\
Entropiebilanz & Maximale innere Entropie & Lokales Minimum via externe Dissipation \\
Erhaltung & Energetische Isolation & Kontinuierlicher metabolischer Umsatz \\
Evolution\"are Rolle & Endpunkt (W\"armetod) & Ausgangspunkt f\"ur Selektion \\
Beispiele & Kristallisiertes Salz & Autokatalytische Zyklen, Zellen \\
\bottomrule
\end{tabularx}
\end{center}

Pr\"abiotische Replikationssysteme unterliegen den DKS-Prinzipien. Die Aufrechterhaltung eines dynamisch-kinetisch stabilen Zustands erfordert unvermeidlich Arbeit, und Arbeit erzeugt unvermeidlich Entropie, die aus dem System exportiert werden muss.

\subsection{Dissipative Strukturen und die entropische Randbedingung}

Die Formulierung von DKS baut tiefgreifend auf Prigogines Arbeit \"uber dissipative Strukturen auf \citep{Bioenergetics}. Dissipative Strukturen sind makroskopisch geordnete Raum-Zeit-Muster, die spontan in Systemen entstehen, die durch starke \"au\ss{}ere Gradienten weit vom Gleichgewicht getrieben werden.

FST-II erweitert diesen Ansatz durch die Erkenntnis, dass pr\"abiotische Chemie letztlich die Suche nach molekularen Netzwerken ist, die als mikroskopische dissipative Strukturen funktionieren. Ein autokatalytisches Reaktionsnetzwerk, das Umgebungsmolek\"ule absorbiert, sie unter Energiefreisetzung umwandelt und Kopien von sich selbst produziert, ist thermodynamisch eine Maschine zur Energiedissipation \citep{Bioenergetics}. Die Stabilit\"at dieses Netzwerks (seine DKS) h\"angt direkt davon ab, wie effizient es diese entropische Randbedingung erf\"ullt.

% ===================================================================
\section{Das Maximum Entropy Production Principle}
\label{sec:mepp}
% ===================================================================

Ein heuristisches Prinzip, das h\"aufig zur Interpretation hochgeordneter DKS-Zust\"ande herangezogen wird, ist das Maximum Entropy Production Principle (MEPP) \citep{EntropyMEPP,MEPP2006}.

\subsection{Formulierung und historische Entwicklung}

In seiner starken Form postuliert MEPP, dass komplexe physikalische und chemische Systeme, die weit vom thermodynamischen Gleichgewicht operieren und \"uber ausreichende Freiheitsgrade verf\"ugen, eine statistische Tendenz zu dynamischen Pfaden und strukturellen Zust\"anden mit hoher Entropieproduktion aufweisen \citep{ThermodynamicsOrganisms}. Diese Arbeit setzt die starke Form nicht als allgemein g\"ultig voraus; sie verwendet MEPP als Kandidaten f\"ur einen Stabilit\"atsfilter, dessen chemische Relevanz unabh\"angig getestet werden muss.

Rod Swenson zeigte, dass die Natur geordnete Strukturen nicht trotz, sondern gerade wegen des Zweiten Hauptsatzes produziert \citep{ThermodynamicsLife}. Geordnete Systeme (wie ein Tornado oder eine lebende Zelle) sind bei weitem effizienter darin, Energiegradienten abzubauen, als ungeordnete Systeme. Die Entstehung des Lebens ist unter dieser Sicht ein physikalisch erwarteter Phasen\"ubergang der planet\"aren Geochemie, um den massiven solaren Energiegradienten effizienter abzubauen.

\subsection{Die England-Ungleichung}

Jeremy England lieferte einen Meilensteinbeitrag, indem er eine untere Schranke f\"ur die Entropieproduktion bei Selbstreplikation aus dem Crooks-Fluktuationstheorem ableitete \citep{Crooks1999,England2013}. Die Herleitung verl\"auft wie folgt.

Betrachte ein isothermes System, das an ein W\"armebad bei Temperatur $T$ gekoppelt ist, mit einem Vorw\"arts- und zeitumgekehrten Protokoll. Sei $A \to B$ ein grobk\"orniger Makro\"ubergang (z.\,B.\ ein erfolgreiches Replikationsereignis), mit Vorw\"artswahrscheinlichkeit $p_{\mathrm{grow}} := p_F(A \to B)$ und R\"uckw\"artswahrscheinlichkeit $p_{\mathrm{decay}} := p_R(B \to A)$. Crooks' Fluktuationstheorem verkn\"upft die Wahrscheinlichkeiten mikroskopischer Vorw\"arts- und R\"uckw\"artsbahnen mit der exponentiierten Entropieproduktion. Nach der Grobk\"ornung von Mikrobahnen zum Makro\"ubergang und Anwendung der Jensenschen Ungleichung ist die relevante Konsequenz die folgende untere Schranke. Wir formulieren diese Jensen-Folge direkt, weil die Vorzeichenkonvention des Exponentialmittels davon abh\"angt, ob W\"armeabgabe an das Bad oder am System geleistete Arbeit parametrisiert wird; die einfache Schreibweise $p_{\mathrm{grow}}/p_{\mathrm{decay}}=\langle e^{+\beta W_{\mathrm{diss}}}\rangle$ w\"urde die Richtung der Jensen-Ungleichung umkehren.

\begin{lemma}[England (2013), aus Crooks FT]
\label{lem:england}
Die mittlere dissipierte Arbeit bedingt auf einen Makro\"ubergang erf\"ullt
\begin{equation}
\langle W_{\mathrm{diss}} \rangle_{A \to B} \geq k_B T \ln \frac{p_{\mathrm{grow}}}{p_{\mathrm{decay}}}.
\end{equation}
\"Aquivalent in Entropieform:
\begin{equation}
\Delta S_{\mathrm{ext}} \geq k_B \ln \frac{p_{\mathrm{grow}}}{p_{\mathrm{decay}}} + \frac{\Delta F_{\mathrm{int}}}{T},
\end{equation}
wobei $\Delta F_{\mathrm{int}}$ die interne Freie-Energie-\"Anderung ist, die mit dem Makro\"ubergang assoziiert ist.
\end{lemma}

\paragraph{Raten-Approximation.} Wenn Wachstum und Zerfall als seltene Ereignisse mit n\"aherungsweise konstanten Hazard-Raten $r_{\mathrm{rep}}$ und $r_{\mathrm{decay}}$ modelliert werden, reduziert sich das Wahrscheinlichkeitsverh\"altnis auf $p_{\mathrm{grow}} / p_{\mathrm{decay}} \approx r_{\mathrm{rep}} / r_{\mathrm{decay}} = \tau_{\mathrm{decay}} / \tau_{\mathrm{rep}}$, was die in der Literatur verwendete heuristische Form $\Delta S_{\mathrm{ext}} + \Delta S_{\mathrm{int}} \geq k_B \ln(\tau_{\mathrm{decay}} / \tau_{\mathrm{rep}})$ wiedergibt. Diese Approximation verwirft die vollst\"andige Pfadensemble-Struktur des Crooks-Rahmens.

\paragraph{Spieltheoretische Lesart.} Die Ungleichung schr\"ankt erreichbare Wachstumsvorteile ein: die Erh\"ohung des Wachstumsproxys $\ln(p_{\mathrm{grow}} / p_{\mathrm{decay}})$ erfordert mindestens proportionale dissipierte Arbeit. Die Thermodynamik erzwingt somit eine \emph{Eintrittskostenbeschr\"ankung} f\"ur den Bereich erreichbarer Wachstumsvorteile---h\"ohere replikative Fitness verlangt entsprechend h\"ohere Entropieproduktion.

\textbf{Vorsicht bzgl.\ Kausalrichtung:} Eine h\"aufig anzutreffende Fehldarstellung von Englands Ergebnis besagt, dass hohe Entropieproduktion Replikation \emph{antreibt} oder \emph{verursacht}. Dies kehrt die tats\"achliche kausale Struktur um. Englands Ungleichung zeigt, dass \emph{erfolgreiche} Replikatoren (hohe Replikationsrate $1/\tau_{rep}$, lange Lebensdauer $\tau_{decay}$) \emph{notwendigerweise mehr Entropie produzieren}---Entropieproduktion ist eine \emph{Konsequenz} erfolgreicher Replikation, nicht ihre Ursache. Die MEPP-als-Selektionsprinzip-Lesart (hohe Entropieproduktion \emph{selektiert f\"ur} Replikatoren) ist eine zus\"atzliche, umstrittene Inferenz, die \"uber Englands eigene Aussagen hinausgeht \citep{England2013}. England selbst war vorsichtig bez\"uglich starker teleologischer Interpretationen. Diese Arbeit verwendet MEPP als heuristisches Organisationsprinzip, w\"ahrend sie anerkennt, dass die kausale Richtung von Entropieproduktion zu Selektion noch zu etablieren ist.

\begin{remark}[Englands Schranke als Stabilit\"atsfilter]
\label{rem:england-stability-filter}
Die Resolventenanalyse aus dem eichtheoretischen Riemannschen-Vermutungs-Programm \citep{GeigerRH2026b} liefert eine strukturelle Kl\"arung von Englands Ungleichung. Die Dissipationsschranke ist kein kausaler Treiber der Replikation, sondern ein \emph{Stabilit\"atsfilter}: replizierende Systeme, die die Schranke verletzen, sind resolventeninstabil und werden durch die Dynamik eliminiert. Unter den vielen molekularen Konfigurationen, die die Schranke in f\"uhrender Ordnung erf\"ullen (die entartete Mannigfaltigkeit), erfolgt die Selektion in zweiter Ordnung durch Kopplung zwischen den replikativen und dissipativen Freiheitsgraden. Diese Resolventenperspektive reduziert die kausale Ambiguit\"at: Englands Ungleichung markiert die \emph{Grenze der lebensf\"ahigen Mannigfaltigkeit}, w\"ahrend MEPP-als-Stabilit\"atsfilter einen resolventenstabilen Kandidaten innerhalb dieser identifiziert. Systeme werden nicht in Richtung hoher Entropieproduktion getrieben; vielmehr persistieren nur diejenigen Konfigurationen, die unter der Dissipationseinschr\"ankung in zweiter Ordnung stabil sind.
\end{remark}

\subsection{\"Uberwindung der Teleologie}

Diese Einsichten ver\"andern das Verst\"andnis pr\"abiotischer Evolution \citep{UnbecomingBeing}. Ein physikalisches System, das durch stochastische thermische Fluktuationen eine Netzwerkarchitektur erreicht, die Replikation und Dissipation zugleich zul\"asst, kann innerhalb seiner Population statistisch persistenter werden, wenn der entsprechende Fixpunkt dynamisch stabil ist.

Im Geiste der Argumente von Sean Carroll \citep{UnbecomingBeing} und Jeremy England \citep{England2013} kann die scheinbare Zielgerichtetheit des Lebens metabolisch umbeschrieben werden: ein Organismus oder ein pr\"abiotischer Hyperzyklus ist eine Struktur, die Replikation mit Energiedissipation koppelt. ``Ein gro\ss{}artiger Weg, mehr Energie zu dissipieren, besteht darin, mehr Kopien von sich selbst zu machen'' \citep{England2013}. Leben muss damit nicht aus einem inh\"arenten \"Uberlebensinstinkt nach Reproduktion streben; Reproduktion kann als thermodynamisch konsistentes Ergebnis behandelt werden, sofern MEPP als eingeschr\"ankte Stabilit\"atsaussage und nicht als aktiver kausaler Treiber gelesen wird.

% ===================================================================
\section{Chemische Spieltheorie: Molek\"ule als Strategietr\"ager}
\label{sec:cgt}
% ===================================================================

Die FST-II-Theorie \"ubersetzt die globalen thermodynamischen Randbedingungen von MEPP in die formale, mikro- und mesoskopische Sprache der Spieltheorie, konkret in das aufstrebende Feld der Chemical Game Theory (CGT) \citep{CGT2018,CGTMacmillan}.

\textbf{Hinweis zur Richtung der CGT-Metapher.} Velegol et al.~\citep{CGT2018} f\"uhrten CGT als Rahmen ein, in dem \emph{menschliche strategische Entscheidungen} mit dem Formalismus chemischer Reaktionen modelliert werden---Spielerentscheidungen werden als metaphorische ``Molek\"ule'' (``Knowlecules'' genannt) behandelt, und Gleichgewichtsergebnisse werden mit chemischer Thermodynamik berechnet. Anders ausgedr\"uckt wendet Velegols CGT Chemie \emph{auf} soziale Spiele an, nicht Spieltheorie \emph{auf} Chemie. Der vorliegende Artikel kehrt diese Metaphernrichtung um: wir wenden spieltheoretische Konzepte \emph{auf} tats\"achliche chemische Systeme an und behandeln reale molekulare Spezies als Tr\"ager strategieartiger Freiheitsgrade. Diese Umkehrung ist unabh\"angig motiviert durch die experimentelle Arbeit von Yeates et al.~\citep{AzoarcusPNAS}, die zeigten, dass pr\"abiotische RNA-Replikator-Dynamiken direkt auf klassische $2\times 2$-Payoff-Matrizen abgebildet werden k\"onnen. Wir behalten die Bezeichnung ``CGT'' f\"ur die Kontinuit\"at mit der Literatur bei, aber der Leser sollte sich bewusst sein, dass unsere Anwendung auf die pr\"abiotische Chemie eine Erweiterung von---nicht eine direkte Anwendung von---Velegols originalem Rahmen darstellt.

\subsection{Spieltheoretische Formulierung}

In der FST-II-Formulierung werden Spieler durch chemische Spezies, Molek\"ule oder funktionale Oligomere repr\"asentiert. Ihre ``Entscheidungen'' oder ``Strategien'' entsprechen probabilistischen Reaktionswahrscheinlichkeiten, sterischen Konformations\"anderungen und spezifischen Katalysepr\"aferenzen, streng bestimmt durch molekulare Erkennungssequenzen und quantenchemische Eigenschaften \citep{ThermodynamicSelection}.

\begin{definition}[Chemisches Spiel]
Ein chemisches Spiel ist ein Tupel $\mathcal{G} = (\mathcal{S}, \mathcal{R}, \Pi, \mathcal{C})$ wobei:
\begin{itemize}
    \item $\mathcal{S} = \{X_1, \ldots, X_N\}$ die Menge der $N$ molekularen Spezies (Spieler) ist
    \item $\mathcal{R} = \{R_1, \ldots, R_m\}$ die Menge der m\"oglichen Reaktionen (Strategieraum) ist
    \item $\Pi: \mathcal{R} \to \mathbb{R}$ die Payoff-Funktion ist
    \item $\mathcal{C} = (T, p, \{c_i\}, \ldots)$ die thermodynamischen Einschr\"ankungen sind
\end{itemize}
\end{definition}

\subsection{Payoff-Funktionen und Gibbs-Energie}

Die Payoff-Funktion jedes Spielers ist definiert als die Minimierung der Gibbs-freien Energie ($\Delta G$) der spezifischen Reaktion, die er katalysiert, unter der starren Einschr\"ankung elementarer Massenbilanz \citep{MetabolicNash}. F\"ur ein komplexes Reaktionsgemisch aus $N$ verschiedenen molekularen Verbindungen und $m$ elementaren Bausteinen wird das Minimierungsproblem formalisiert als:

\begin{equation}
\min_{\{n_i\}} G = \sum_{i=1}^{N} n_i \left( G_i^0 + RT \ln\frac{n_i}{n_{\mathrm{total}}} \right)
\end{equation}

\noindent
wobei $n_i$ die Stoffmenge der Spezies $i$ und $n_{\mathrm{total}} = \sum_i n_i$ die Gesamtstoffmenge bezeichnet.

\noindent
\textbf{Idealit\"atsannahme und Aktivit\"atskoeffizienten:} Dieser Ausdruck verwendet die ideale Mischungsentropie $RT\ln(n_i/n_{\mathrm{total}})$, die streng nur f\"ur ideale (verd\"unnte) L\"osungen gilt. In konzentrierten L\"osungen m\"ussen Konzentrationen durch Aktivit\"aten $a_i = \gamma_i c_i$ ersetzt werden, wobei $\gamma_i$ der Aktivit\"atskoeffizient der Spezies $i$ ist; wir arbeiten durchgehend im verd\"unnten Grenzfall $\gamma_i \approx 1$. Pr\"abiotische Mischungen mit starken intermolekularen Wechselwirkungen (autokatalytische Kopplung, Kreuzinhibition) erfordern diese Aktivit\"atskoeffizienten, die $\ln(n_i/n_{\mathrm{total}})$ durch $\ln(\gamma_i n_i/n_{\mathrm{total}})$ ersetzen. Die ideale Form wird hier als N\"aherung f\"uhrender Ordnung beibehalten; die qualitative spieltheoretische Struktur (Stationarit\"atsbedingungen, Nash-Gleichgewichts-Analogie) bleibt unter nicht-idealen Korrekturen erhalten, obwohl sich die spezifischen Gleichgewichtszusammensetzungen verschieben.

Eine notwendige Bedingung f\"ur ein Nash-Gleichgewicht in diesem biochemischen Kontext ist, dass kein einzelner molekularer Reaktionspfad (Spieler) seine spezifische Gibbs-Energie durch eine einseitige \"Anderung seiner katalytischen Rate oder seines Flusses weiter minimieren kann, ohne globale Massenbalance-Restriktionen zu verletzen \citep{MetabolicNash}. \textbf{Wichtige Einschr\"ankung:} Diese Stationarit\"atsbedingung erster Ordnung ist notwendig, aber nicht hinreichend f\"ur ein Nash-Gleichgewicht im vollen spieltheoretischen Sinne. Die Gibbs-Minimierung liefert ein eindeutiges globales Minimum (konvexe Optimierung), w\"ahrend Nash-Gleichgewichte nicht-eindeutig sein k\"onnen und Sattelpunkte einschlie\ss{}en. Was hier etabliert wird, ist eine formale Analogie---die station\"aren Punkte der eingeschr\"ankten Gibbs-Minimierung erf\"ullen die notwendigen Bedingungen f\"ur ein Nash-Gleichgewicht---kein vollst\"andiger Isomorphismus.

\subsection{Verbindung zu MEPP}

In isothermen isobaren \emph{geschlossenen} Systemen ist die \"Anderung der Gibbs-freien Energie ($\Delta G$) direkt mit der totalen Entropieproduktion ($\Delta S_{total}$) verkn\"upft durch $\Delta G = -T\Delta S_{total}$. Diese Identit\"at gilt streng nur f\"ur geschlossene Systeme bei konstantem Druck und konstanter Temperatur. F\"ur offene NESS-Systeme (nichtgleichgewichtige station\"are Zust\"ande) muss der Gibbs-Freie-Energie-Ansatz um explizite Entropieaustauschterme mit der Umgebung erg\"anzt werden; die Relation gilt n\"aherungsweise, wenn Stofffl\"usse im Vergleich zu Reaktionsraten langsam sind. In diesem eingeschr\"ankten Setting kann ein molekulares System, das $\Delta G$ schnell und stark minimiert, zugleich die Entropieproduktion erh\"ohen. Ein tiefes Nash-Gleichgewicht mit hohem aggregiertem Payoff im chemischen Netzwerk zu finden, ist daher als Identifikation einer Kandidatenkonfiguration unter MEPP- und England-artigen Randbedingungen zu verstehen, nicht als Beweis daf\"ur, dass Entropieproduktion diese Konfiguration kausal ausw\"ahlt.

\textbf{Zirkularit\"ats-Warnung:} Die Argumentstruktur hier riskiert Zirkularit\"at: (1) MEPP wird als globale Payoff-Funktion verwendet, (2) Nash-Gleichgewichte werden als Maxima dieses MEPP-Payoffs definiert, (3) molekulare Systeme, die Nash-Gleichgewichte erreichen, werden als MEPP-maximierend vorhergesagt. Dies ist eine Definitionsschleife, kein unabh\"angiger empirischer Test. Der Rahmen w\"are nur dann nicht-zirkul\"ar, wenn MEPP und Nash-Gleichgewichte unabh\"angig definiert w\"aren und ihr gemeinsames Auftreten empirisch getestet werden k\"onnte. Abschnitt~\ref{sec:predictions} (P3) liefert einen Test dieser Art---die Messung der Entropieproduktion konkurrierender autokatalytischer Netzwerke unabh\"angig von ihrem spieltheoretischen Status---aber dieser Test wurde noch nicht durchgef\"uhrt. Bis solche Tests berichtet werden, sollte die MEPP-als-Payoff-Funktion-Behauptung als Arbeitshypothese mit heuristischem Wert behandelt werden, nicht als etabliertes Ergebnis.

% ===================================================================
\section{Empirische Validierung am Azoarcus-Ribozym}
\label{sec:azoarcus}
% ===================================================================

Die Postulate der Chemical Game Theory werden empirisch motiviert durch die Arbeit an Azoarcus-Ribozymen von Niles Lehman und Kollegen \citep{AzoarcusPNAS,AzoarcusPMC}. Das Azoarcus-System ist ein besonders relevantes Modell, weil es innerhalb der RNA-Welt-Hypothese operiert---dem breit diskutierten Szenario, dass selbstreplizierende RNA-Molek\"ule dem DNA-Protein-Leben vorausgingen \citep{Gilbert1986,Eigen1979}---und eine der wenigen experimentellen Plattformen bietet, auf denen pr\"abiotische Replikator-Dynamiken direkt beobachtet werden k\"onnen. \textbf{Einschr\"ankung des G\"ultigkeitsbereichs:} Das Azoarcus-System liefert 16 Genotypen (aus einem viel gr\"o\ss{}eren theoretischen Sequenzraum) in einem sorgf\"altig kontrollierten In-vitro-Experiment unter spezifischen Magnesium- und Temperaturbedingungen. Die Extrapolation von diesem einzelnen experimentellen System auf universelle Prinzipien der pr\"abiotischen Chemie ist eine induktive Inferenz, die vorsichtig gemacht werden sollte. Replikation der spieltheoretischen Klassifikation \"uber verschiedene RNA-Systeme, verschiedene Molek\"ulklassen (Peptide, Lipide) und verschiedene Umweltbedingungen ist erforderlich, bevor der Rahmen Allgemeing\"ultigkeit beanspruchen kann.

\subsection{Das experimentelle System}

Das Azoarcus-Ribozym ist ein katalytisch aktives RNA-Molek\"ul von ungef\"ahr 200 Nukleotiden. Es kann in Fragmente (WXY und Z, jeweils $\sim$50 Nukleotide) aufgespalten werden, die in warmer Magnesiuml\"osung spontan und kovalent zu intakten, autokatalytisch aktiven WXYZ-Molek\"ulen assemblieren \citep{AzoarcusPMC}:
\begin{equation}
\text{WXY} + \text{Z} + \text{WXYZ} \longrightarrow 2\,\text{WXYZ}
\end{equation}

Die Spezifit\"at der molekularen Assemblierung wird durch molekulare Erkennung basierend auf 3-Nukleotid-Basenpaarungen zwischen Internal Guide Sequence (IGS) und komplement\"aren ``Tag''-Sequenzen kontrolliert \citep{AzoarcusPNAS}. Dieses System baut auf fr\"uheren Demonstrationen template-gesteuerter Selbstreplikation in nukleotidbasierten Oligomeren durch Sievers \& von Kiedrowski \citep{SieversKiedrowski1994} auf, die zeigten, dass die kreuzkatalytische Replikation komplement\"arer Templates mit Effizienzen verlaufen kann, die mit autokatalytischer Selbstreplikation vergleichbar sind.

\subsection{Die vier chemischen Spiele}

Durch Mutation des mittleren Nukleotids von IGS und Tag erzeugten die Forscher 16 m\"ogliche molekulare Genotypen. In Konkurrenzszenarien zwischen zwei RNA-Genotypen k\"onnen die Interaktionen auf klassische 2$\times$2-Payoff-Matrizen abgebildet werden \citep{AzoarcusPMC}:

\begin{center}
\setlength{\tabcolsep}{4pt}
\begin{tabularx}{\textwidth}{@{}>{\raggedright\arraybackslash}p{0.18\textwidth}>{\raggedright\arraybackslash}p{0.28\textwidth}>{\raggedright\arraybackslash}X>{\raggedright\arraybackslash}p{0.12\textwidth}@{}}
\toprule
\textbf{Spieltyp} & \textbf{Kinetische Bedingung} & \textbf{Molekulare Dynamik} & \textbf{Beispiel} \\
\midrule
Dominanz & $k_{AA} > k_{BA}$, $k_{AB} > k_{BB}$ & Eine Spezies verdr\"angt alle & CG vs.\ GA \\
Kooperation & $k_{BA} > k_{AA}$, $k_{AB} > k_{BB}$ & Gegenseitige Kreuzkatalyse dominiert & AC vs.\ UU \\
Hirschjagd & $k_{AA} > k_{BA}$, $k_{BB} > k_{AB}$, $k_{AA} > k_{BB}$ & Frequenzabh\"angige Bistabilit\"at & AU vs.\ UC \\
Gefangenen-\\dilemma & $k_{BA} > k_{AA} > k_{BB} > k_{AB}$ & Parasit beutet Kooperator aus & CA vs.\ GG \\
\bottomrule
\end{tabularx}
\end{center}

\noindent
Bei der Hirschjagd sind beide reinen Strategien (all-A und all-B) Nash-Gleichgewichte, aber das payoff-dominante Gleichgewicht (all-A, da $k_{AA} > k_{BB}$) erfordert koordinierte Adoption und ist daher frequenzabh\"angig: jede Spezies gewinnt nur, wenn sie bereits in der Mehrheit ist.

Das Gefangenendilemma erfordert die vollst\"andige Ordnung $T > R > P > S$ (in kinetischen Termen: $k_{BA} > k_{AA} > k_{BB} > k_{AB}$) zusammen mit der Stabilit\"atsbedingung des iterierten Spiels $2R > T + S$ (d.\,h.\ $2k_{AA} > k_{BA} + k_{AB}$), die verhindert, dass alternierende Ausbeutung die gegenseitige Kooperation \"ubertrifft \citep{HofbauerSigmund1998}.

\subsection{Das molekulare Gefangenendilemma}

Das empirisch demonstrierte molekulare Gefangenendilemma enth\"ullt ein tiefes konzeptuelles Problem: ein lokaler thermodynamischer Vorteil eines isolierten egoistischen Molek\"uls (des Parasiten) kann zum Zusammenbruch des gesamten kooperativen Netzwerks f\"uhren \citep{AzoarcusPMC}. Dies w\"urde das System beenden und die Entropieproduktion abbrechen.

Die Verhinderung dieses parasit\"aren Kollapses bei gleichzeitiger Aufrechterhaltung anhaltender Dissipation erfordert einen Stabilit\"atsmechanismus auf einer h\"oheren Integrationsebene. Dies f\"uhrt mathematisch zur Architektur der Eigen-Schuster-Hyperzyklen.

% ===================================================================
\section{Hyperzyklen, Replikator-Dynamik und r\"aumliche ESS}
\label{sec:hypercycles}
% ===================================================================

\subsection{Die Hyperzyklus-Architektur}

In einem Hyperzyklus bilden $N$ kurze, relativ fehlertolerante Quasispezies $I_1, \ldots, I_N$ einen Ring \citep{Eigen1979}. Jede Spezies katalysiert nicht nur ihre eigene Replikation, sondern f\"ordert auch die Replikation der nachfolgenden Spezies ($I_i$ katalysiert $I_{i+1}$).

\subsection{Die Replikatorgleichung}

FST-II verwendet die Replikatorgleichung als rigorose Formalisierung der Hyperzyklus-Dynamik \citep{ReplicatorEntropy,ReplicatorDiffusion}:

\begin{equation}
\dot{x}_i = x_i \left[ (\mathbf{A}\mathbf{x})_i - \mathbf{x}^T \mathbf{A} \mathbf{x} \right]
\end{equation}

wobei:
\begin{itemize}
    \item $\mathbf{x}$ der Zustandsvektor relativer H\"aufigkeiten auf dem Einheitssimplex ist
    \item $\mathbf{A}$ die Payoff-Matrix der Interaktionen (katalytische Raten) ist
    \item $(\mathbf{A}\mathbf{x})_i$ der erwartete Payoff (Fitness) der Spezies $i$ ist
    \item $\mathbf{x}^T \mathbf{A} \mathbf{x}$ die mittlere Fitness der Population ist
\end{itemize}

Eine molekulare Spezies w\"achst relativ zum Rest der Population genau dann, wenn ihre netzwerkspezifische Fitness die durchschnittliche Fitness \"ubersteigt.

\subsection{Abbildung auf Nash-Gleichgewichte und ESS}

Es existiert eine strenge formale Abbildung zwischen der Asymptotik der Replikatorgleichung und Nash-Gleichgewichten \citep{HofbauerSigmund1998,HofbauerSigmund2003}:

\begin{theorem}[Replikator-Nash-Korrespondenz {\citep{HofbauerSigmund2003,HofbauerSigmund1998}}]
\label{thm:replicator-nash}
Jedes Nash-Gleichgewicht $\mathbf{x}^*$ des durch Matrix $\mathbf{A}$ definierten Spiels ist ein Ruhepunkt der Replikatorgleichung. Umgekehrt ist jeder innere $\omega$-Limespunkt der Replikator-Dynamik ein Nash-Gleichgewicht, und jeder Ljapunow-stabile Ruhepunkt im Inneren des Simplex ist ein Nash-Gleichgewicht. Ferner entspricht jede ESS einem asymptotisch stabilen Ruhepunkt der Replikator-Dynamik.
\end{theorem}

\begin{remark}[Geltungsbereich der Korrespondenz]
Die \"Aquivalenz zwischen inneren Ruhepunkten der Replikator-Dynamik und Nash-Gleichgewichten ist klassisch (Hofbauer \& Sigmund 1998, Theorem~7.2.1). Jedes Nash-Gleichgewicht, auch ein Rand-Nash-Gleichgewicht, ist ein Ruhepunkt der Replikator-Dynamik. Die heikle Richtung ist die Umkehrung: Rand-Ruhepunkte---bei denen eine oder mehrere Spezies ausgestorben sind---m\"ussen nicht Nash-Gleichgewichten entsprechen, weil eine fehlende Spezies bei Wiedereinbringung einen h\"oheren Payoff haben kann. Diese Unterscheidung ist f\"ur die nachfolgenden pr\"abiotischen Anwendungen wichtig, da das Aussterben molekularer Spezies ein physikalisch relevantes Szenario darstellt und Invasionsbedingungen explizit gepr\"uft werden m\"ussen.
\end{remark}

Dieses Ergebnis ist ein Eckpfeiler der klassischen evolution\"aren Spieltheorie \citep{Weibull1995}; der Beitrag des vorliegenden Artikels ist nicht das Theorem selbst, sondern seine Anwendung auf pr\"abiotische chemische Netzwerke und seine Integration mit der MEPP-Stabilit\"atsfilter-Interpretation. Siehe Hofbauer \& Sigmund (1998, 2003) f\"ur Beweise. Eine Evolutionarily Stable Strategy (ESS) ist eine molekulare Populationsverteilung $\mathbf{x}^*$, die gegen Invasion durch jede seltene ``Mutanten''-Strategie $\mathbf{y}$ resistent ist. Eine ESS ist immer ein Nash-Gleichgewicht und impliziert asymptotische Stabilit\"at in der Replikator-Dynamik \citep{HofbauerSigmund1998}.

\begin{remark}[Autokatalytische Stabilit\"at als Resolventenselektion]
\label{rem:autocatalysis-resolvent}
Der Resolventenrahmen aus der Riemannschen-Vermutungs-Analyse \citep{GeigerRH2026b} beleuchtet, warum spezifische autokatalytische Netzwerke (Hyperzyklen, Spiralwellen) persistieren, w\"ahrend andere zerfallen. Diese Strukturen sind keine \emph{globalen} MEPP-Maxima---sie maximieren nicht global die Entropieproduktion \"uber alle denkbaren Konfigurationen. Vielmehr sind sie \emph{in zweiter Ordnung resolventenstabilisierte Netzwerke}: lokale Nash-Gleichgewichte, die unter St\"orung stabil sind. In f\"uhrender Ordnung produzieren viele Netzwerktopologien vergleichbare Replikationsraten und erf\"ullen die England-Dissipationsschranke. Parasit\"are Invasion, r\"aumliche Instabilit\"at und Konzentrationskollaps sind Destabilisierungsmechanismen zweiter Ordnung, die alle au\ss{}er der resolventenstabilen Konfiguration eliminieren. Die Replikatorgleichung selektiert somit den resolventenstabilen Fixpunkt, nicht notwendig den produktivsten. Dies kl\"art die Rolle von MEPP in diesem Rahmen: MEPP dient als notwendige Bedingung (Stabilit\"atsfilter) und nicht als hinreichendes Selektionsprinzip (Maximierung). Reale pr\"abiotische Netzwerke weisen Robustheit auf, weil sie resolventenstabil sind, was mit---aber nicht identisch ist mit---hoher Entropieproduktion korreliert.
\end{remark}

\subsection{Ljapunow-Funktionen und thermodynamische Analogie}

F\"ur Systeme mit symmetrischer Interaktionsmatrix $\mathbf{A}$ (oder, pr\"aziser, wenn das Spiel eine Potentialfunktion zul\"asst) steigt die mittlere Fitness $\phi(\mathbf{x}) = \mathbf{x}^T \mathbf{A} \mathbf{x}$ monoton entlang von Replikator-Trajektorien, bis ein lokales Maximum erreicht wird \citep{ReplicatorEntropy}. Dies ist ein Analogon von Fishers Fundamentaltheorem, adaptiert auf Replikator-Dynamik; das Originaltheorem betrifft additive genetische Varianz in der Populationsgenetik, und seine direkte Anwendbarkeit auf chemische Replikatorsysteme erfordert die Symmetrie- (bzw.\ Potentialspiel-)Bedingung an $\mathbf{A}$. Unter dieser Bedingung fungiert die mittlere Fitness mathematisch als Ljapunow-Funktion.

\paragraph{Entropiedynamik unter Replikatorfluss.}
F\"ur die Replikator-Dynamik $\dot{x}_i = x_i(f_i(x) - \phi(x))$ mit $\phi(x) = \sum_i x_i f_i(x)$ erf\"ullt die Shannon-Entropie $H(x) = -\sum_i x_i \ln x_i$ die exakte Identit\"at:
\begin{equation}
\dot{H}(x) = -\sum_i x_i \big(f_i(x) - \phi(x)\big) \ln x_i = -\mathrm{Cov}_x(f, \ln x),
\end{equation}
die \emph{kein bestimmtes Vorzeichen} hat. Die Shannon-Entropie liefert kein universelles Relaxationsgesetz unter Replikator-Dynamik.

\paragraph{KL-Divergenz als Ljapunow-Funktional.}
Ein thermodynamisch sinnvolles Ljapunow-Funktional ist stattdessen die Kullback-Leibler-Divergenz zu einem Gleichgewichtszustand $\mathbf{x}^*$:
\begin{equation}
D_{\mathrm{KL}}(\mathbf{x}^* \| \mathbf{x}) = \sum_i x_i^* \ln \frac{x_i^*}{x_i},
\end{equation}
deren Zeitableitung entlang Replikator-Trajektorien ist:
\begin{equation}
\frac{d}{dt} D_{\mathrm{KL}}(\mathbf{x}^* \| \mathbf{x}) = -\sum_i x_i^* \big(f_i(x) - \phi(x)\big).
\end{equation}
Unter spezifischen Stabilit\"atsannahmen an $\mathbf{x}^*$---n\"amlich dass $\mathbf{x}^*$ eine innere ESS ist, oder dass das Spiel eine Potentialfunktion zul\"asst (\"aquivalent: $\mathbf{A}$ ist symmetrisch bis auf eine additive Spaltenkonstante)---ist diese Gr\"o\ss{}e nicht-zunehmend und liefert eine $H$-Theorem-artige Aussage: die Populationsverteilung relaxiert zum Nash-Gleichgewicht, w\"ahrend die KL-Divergenz monoton abnimmt \citep{HofbauerSigmund1998,ReplicatorEntropy}. F\"ur allgemeine asymmetrische Payoff-Matrizen muss die KL-Divergenz nicht monoton abnehmen, und die Ljapunow-Eigenschaft versagt.

Dies liefert eine strukturelle Analogie zwischen populationsdynamischer Evolution zu Nash-Gleichgewichten und thermodynamischer Relaxation: die KL-Divergenz spielt die Rolle eines Freie-Energie-Funktionals, und der Replikatorfluss wirkt als Gradientenabstieg in der Shahshahani-Metrik der Informationsgeometrie \citep{ReplicatorEntropy}. Die Verbindung zur thermodynamischen Entropieproduktion (MEPP) erfordert den zus\"atzlichen Schritt, die mittlere Fitness mit der thermodynamischen Dissipationsrate zu identifizieren, was eine Approximation ist, die in spezifischen chemischen Settings g\"ultig ist, aber nicht im Allgemeinen.

\subsection{R\"aumliche Verteilung: Spiralwellen als r\"aumlich verteilte ESS}

Trotz ihrer mathematischen Eleganz leiden gut durchmischte Hyperzyklen an einem fatalen physikalischen Defekt: sie sind extrem anf\"allig f\"ur parasit\"are Molek\"ule---eine molekulare Manifestation der Trag\"odie der Allmende \citep{HypercycleInfinite}.

FST-II integriert r\"aumliche Erweiterung durch Reaktions-Diffusions-Gleichungen \citep{ReplicatorDiffusion}. Wenn Molek\"ule nicht nur reagieren, sondern diffundieren, transformiert sich das Modell in ein System partieller Differentialgleichungen. Das Konzept erweitert sich zu einem \textbf{r\"aumlich verteilten Nash-Gleichgewicht}.

\paragraph{Mean-Field-Gleichgewichts-Formulierung.}
Ein r\"aumlich heterogenes Konzentrationsprofil $\mathbf{x}^*(\mathbf{r})$ stellt ein Mean-Field-Nash-Gleichgewicht dar, wenn es den Fixpunkt-Abschluss erf\"ullt: jede repr\"asentative molekulare Spezies reagiert optimal auf die Populationsverteilung, und die Populationsverteilung ist konsistent mit der induzierten besten Antwort. Konkret, mit $m$ als r\"aumliche Verteilung der Strategien:
\begin{equation}
\mathbf{x}^*(\mathbf{r}) \in \arg\max_{\mathbf{x}} \int_\Omega U(\mathbf{x}, m^*, \mathbf{r}) \, d\mathbf{r}, \qquad m^* = \mathcal{L}(\mathbf{x}^*),
\end{equation}
wobei die Konsistenzbedingung $m^* = \mathcal{L}(\mathbf{x}^*)$ dies von einem globalen (Planer-)Optimum unterscheidet. Im Spezialfall, dass die chemische Interaktion ein Potential $\Phi$ zul\"asst (wie es \"ublich ist f\"ur Reaktionsnetzwerke, die Detailed Balance erf\"ullen, wo die freie Energie als Ljapunow-Funktion dient), fallen Nash-Gleichgewichte mit Maximierern von $\Phi$ zusammen, und die $\arg\max$-Formulierung ist exakt.

Wir f\"uhren den Begriff \textbf{Distributed Evolutionarily Stable Strategy (DESS)} f\"ur eine r\"aumlich heterogene L\"osung ein, die diese Fixpunkt-Bedingung erf\"ullt und zus\"atzlich Stabilit\"at im mittleren Integralsinn aufrechterh\"alt---d.\,h.\ die r\"aumlich gemittelte Fitness des residenten Profils kann durch keine lokalisierte Mutanteninvasion \"ubertroffen werden. Das DESS-Konzept ist \emph{keine} Standardterminologie der evolution\"aren Spieltheorie; es ist unsere Erweiterung der klassischen ESS, die f\"ur gut durchmischte Populationen gilt, auf r\"aumlich verteilte Reaktions-Diffusions-Systeme. Den mathematischen Rahmen f\"ur r\"aumliche Replikator-Diffusions-Dynamik, einschlie\ss{}lich Wanderwellenanalyse konkurrierender ESS, entwickelten Hutson \& Vickers~\citep{ReplicatorDiffusion}; die DESS-Terminologie und ihre Anwendung auf pr\"abiotische Hyperzyklen sind Beitr\"age des vorliegenden Artikels.

R\"aumlich verteilte Replikatorsysteme k\"onnen Homogenit\"at brechen und makroskopische dissipative Strukturen in Form von \textbf{Spiralwellen} bilden \citep{SpiralWaves}. In diesen Spiralwellen rotieren die Populationen autokatalytischer Spezies in konzentrischen Gradienten. Boerlijst \& Hogeweg~\citep{SpiralWaves,HypercycleInfinite} demonstrierten diese r\"aumliche Selbststrukturierung mit zellul\"aren Automaten und Reaktions-Diffusions-Simulationen; sie zeigten, dass egoistische Parasiten, die das System lokal kollabieren lassen, in der Peripherie isoliert bleiben k\"onnen, wo sie an Ressourcenknappheit zugrunde gehen, w\"ahrend der produktive Kern des Hyperzyklus im Zentrum ungest\"ort rotiert. \textbf{Unser Beitrag} ist die Neuinterpretation dieser r\"aumlichen Robustheit als DESS innerhalb des hier entwickelten spieltheoretischen Rahmens---Boerlijst \& Hogeweg verwendeten in ihrer urspr\"unglichen Analyse keine spieltheoretische Sprache und keine entsprechenden Konzepte.

Die Entstehung von Spiralwellen wird hier nicht als zuf\"alliges Muster, sondern als robuste topologische Struktur interpretiert, die eine DESS im Einklang mit den dissipativen Randbedingungen von MEPP aufrechterhalten kann \citep{HypercycleInfinite}.

% ===================================================================
\section{Homochiralit\"at als Symmetriebrechung}
\label{sec:chirality}
% ===================================================================

Die vorangehenden Abschnitte haben argumentiert, dass autokatalytische Netzwerke und Hyperzyklen als Nash-Gleichgewichts-artige Strukturen verstanden werden k\"onnen, die durch r\"aumliche Selbstorganisation gegen parasit\"are Invasion stabilisiert werden. Wir wenden uns nun einem komplement\"aren Ph\"anomen zu---der Chiralit\"atsselektion---das derselbe spieltheoretische Rahmen als eine andere Art von Symmetriebrechung interpretiert: nicht die r\"aumliche Brechung der Translationssymmetrie (Spiralwellen), sondern die Brechung der Spiegelsymmetrie zwischen Enantiomeren.

Eine wichtige Anwendung von FST-II ist eine spieltheoretische Einordnung eines zentralen Problems der pr\"abiotischen Chemie: des Ursprungs der biologischen Homochiralit\"at \citep{Blackmond2010,Biochirality}.

\subsection{Das Chiralit\"ats-R\"atsel}

Leben auf der Erde ist asymmetrisch; es basiert fast ausschlie\ss{}lich auf L-Aminos\"auren (in Proteinen) und D-Zuckern (in DNA/RNA) \citep{HomochiralityThermo}. In der klassischen abiotischen Chemie f\"uhren Synthesen chiraler Molek\"ule ohne asymmetrische Katalysatoren stets zu racemischen Mischungen (50\% L, 50\% D). Ein solches Racemat repr\"asentiert den Zustand maximaler Entropie in einem isolierten System.

\subsection{Das Racemat als Nash-Gleichgewicht gemischter Strategien}

Aus der CGT-Perspektive gleicht das racemische Gemisch exakt einem Spiel mit gemischten Strategien (Mixed Strategy Nash Equilibrium). Beide ``Spieler'' (L-Konfiguration und D-Konfiguration-Molek\"ule) haben identische Payoff-Funktionen in Abwesenheit anderer Einschr\"ankungen und interagieren mit gleicher Wahrscheinlichkeit. In einem geschlossenen System nahe dem Gleichgewicht ist dieser Zustand durch Zeitumkehr-Symmetrie gesch\"utzt und thermodynamisch stabil \citep{PlassonReview2007,TimeReversalSymmetry}.

Jedoch argumentiert FST-II aus der Perspektive stark dissipativer Nicht-Gleichgewichtsthermodynamik (weit vom Gleichgewicht). In einem getriebenen System, das durch konstanten Energiefluss offen gehalten wird, verliert das gemischte Nash-Gleichgewicht des Racemats seine Ljapunow-Stabilit\"at und transformiert sich in einen energetischen Sattelpunkt \citep{Frank1953,PlassonReview2007}.

\subsection{Das Frank-Modell: Autokatalyse und Kreuzinhibition}

\textbf{Priorit\"atshinweis:} Der Mechanismus autokatalytischer Verst\"arkung kombiniert mit Kreuzinhibition, der zu Homochiralit\"at f\"uhrt, wurde von Frank (1953) \citep{Frank1953} eingef\"uhrt, lange vor dem spieltheoretischen Vokabular. Die mathematische Struktur dieser Symmetriebrechung ist in der Chemieliteratur gut etabliert \citep{PlassonReview2007}. Der Beitrag von FST-II ist die Neuinterpretation dieses bekannten Mechanismus innerhalb eines spieltheoretischen Rahmens (\"Ubergang vom gemischten zum reinen Nash-Gleichgewicht) und seine Integration mit MEPP. Ob diese Neuinterpretation pr\"adiktive Vorteile gegen\"uber bestehenden Modellen bietet---oder blo\ss{} eine Umformulierung bekannter Ergebnisse darstellt---ist die Frage, die die testbaren Vorhersagen in Abschnitt~\ref{sec:predictions} (P2) beantworten sollen.

Der formale Mechanismus dieser Destabilisierung beruht auf dem Frank-Modell \citep{Frank1953}:
\begin{enumerate}
    \item \textbf{Autokatalytische Verst\"arkung:} L-Molek\"ule katalysieren bevorzugt die Synthese weiterer L-Molek\"ule aus achiralen Vorl\"aufern (ebenso f\"ur D)
    \item \textbf{Kreuzinhibition:} L- und D-Molek\"ule k\"onnen einen inaktiven Heterodimer-Komplex (L-D) bilden, der aus dem katalytischen Zyklus entfernt wird \citep{Biochirality}
\end{enumerate}

In diesem Regime kann schon eine kleine stochastische Konzentrationsabweichung---etwa ein zuf\"alliger \"Uberschuss von L-Enantiomeren---durch positive Autokatalyse verst\"arkt werden. Gleichzeitig unterdr\"uckt die Kreuzinhibition den bereits benachteiligten D-Antagonisten, bis er im idealisierten Modell verschwindet.

\subsection{Viedma-Ripening: Empirische Unterst\"utzung}

Dass dieser Mechanismus der theoretischen Symmetriebrechung ein experimentelles Analogon hat, wurde durch Viedma-Deracemisierung demonstriert \citep{HomochiralityThermo}. \textbf{Geltungsbereichshinweis:} Viedma-Ripening betrifft Kristallisation unter mechanischem Schleifen, nicht pr\"abiotische autokatalytische Netzwerke. Es liefert empirische Unterst\"utzung f\"ur den \emph{Mechanismus} autokatalytischer Chiralit\"atsverst\"arkung, nicht einen direkten Test des FST-II-Rahmens unter pr\"abiotischen Bedingungen. In Experimenten wurden racemische Suspensionen chiraler Kristalle unter konstantem mechanischem Energieeintrag (Mahlen mit Glasperlen) weit vom thermodynamischen Gleichgewicht im DKS-Regime gehalten.

Das \"uberraschende Ergebnis: die Population rechts- und linksh\"andiger Kristalle koexistiert unter diesen Bedingungen nicht stabil. Durch irreversibles autokatalytisches Aufl\"osen und Wachsen verschwindet eine chirale Population, w\"ahrend die andere sich vollst\"andiger 100\% optischer Reinheit (reine Strategie) ann\"ahert \citep{HomochiralityThermo}.

\subsection{MEPP-Erkl\"arung}

Ein gemischtes, racemisches Polymer- oder Hyperzyklus-Netzwerk leidet unter immenser struktureller Frustration. Falsche Basenpaarungen, sterische Hinderungen und Bildung inaktiver L-D-Heterodimere reduzieren die globale Replikationsrate dramatisch. Eine resolventenstabile Netzwerkkonfiguration---eine Konfiguration, in der Destabilisierungskan\"ale zweiter Ordnung zugleich unterdr\"uckt werden---erfordert topologische Koordination der chemischen Strategien \citep{Biochirality}.

Makroskopische Chiralit\"ats-Symmetriebrechung in der pr\"abiotischen Chemie wird daher nicht als blo\ss{}er historischer Zufall modelliert, sondern als robuste Modellfolge \emph{unter den Bedingungen des Frank-Modells} (Autokatalyse mit gegenseitiger Kreuzinhibition). Homochiralit\"at repr\"asentiert das resolventenstabile Gleichgewicht reiner Strategien f\"ur komplexe pr\"abiotische Netzwerke: innerhalb des Frank-Modells mit zwei Spezies ist es die Konfiguration, bei der die Kreuzkopplungen zweiter Ordnung zwischen chiralen Sektoren gleichzeitig stabilisiert sind \citep{Biochirality}. Ob zus\"atzliche resolventenstabile Konfigurationen in komplexeren Multi-Chiralit\"atszentren-Systemen existieren, bleibt eine offene Frage.

\textbf{Bedingte vs.\ absolute Notwendigkeit.} Die Behauptung der ``Unvermeidlichkeit'' ist auf die Pr\"amissen des Frank-Modells beschr\"ankt---autokatalytische Verst\"arkung und Kreuzinhibition. Ob diese Pr\"amissen selbst in pr\"abiotischen Umgebungen regelm\"a\ss{}ig auftreten, ist eine separate und offene Frage, die FST-II nicht l\"ost. Goulds Kontingenzargument~\citep{Gould1989}---dass ein Abspielen des Bandes des Lebens zu radikal anderen Ergebnissen f\"uhren k\"onnte (siehe auch Carroll~\citep{UnbecomingBeing} f\"ur eine moderne Behandlung)---betrifft genau die Frage, ob diese Voraussetzungen universell entstehen. FST-II st\"utzt die engere Aussage, dass \emph{gegeben} Autokatalyse und Kreuzinhibition Homochiralit\"at ein robust erwartbares Ergebnis ist; es stellt nicht fest, dass die Voraussetzungen selbst notwendig sind. Die Unterscheidung zwischen bedingter und absoluter Notwendigkeit sollte durchgehend beachtet werden.

\begin{remark}[Homochiralit\"at als Resolventeneffekt zweiter Ordnung]
\label{rem:chirality-resolvent}
Die Resolventenanalyse aus dem Riemannschen-Vermutungs-Programm \citep{GeigerRH2026b} liefert eine pr\"azise strukturelle Analogie f\"ur Chiralit\"ats-Symmetriebrechung. Das Racemat ist ein in f\"uhrender Ordnung entartetes Gleichgewicht: L- und D-Enantiomere haben identische thermodynamische Payoffs erster Ordnung, analog zu den geraden und ungeraden Sektoren der Weil-quadratischen Form, die identische Laplace-Limiten teilen. Kein energetischer Vorteil unterscheidet die beiden in f\"uhrender Ordnung. Chiralit\"at entsteht als \emph{Effekt zweiter Ordnung} durch Kopplung zwischen den Reaktionsunterr\"aumen: autokatalytische Verst\"arkung und Kreuzinhibition bilden die Resolventenkopplung, die die entartete Mannigfaltigkeit aufspaltet. Der selektierte homochirale Zustand (reines L oder reines D) ist der resolventenstabile Fixpunkt---nicht weil er in erster Ordnung mehr Entropie produziert als das Racemat, sondern weil er die einzige Konfiguration ist, bei der die Kreuzkopplungen zweiter Ordnung zwischen chiralen Sektoren gleichzeitig stabilisiert sind.
\end{remark}

% ===================================================================
\section{Synthese: Die hierarchische Kaskade}
\label{sec:synthesis}
% ===================================================================

Die Architektur der pr\"abiotischen Evolution kann innerhalb von FST-II als dreistufige statistisch gerichtete Kaskade beschrieben werden \citep{ThermodynamicsLife}:

\subsection{Mikroebene: Thermodynamik der Knoten}

Auf der niedrigsten Ebene werden Molek\"ule und Enzyme als lokale Reaktionskan\"ale modelliert. Sie folgen der klassischen Gleichgewichtsthermodynamik, wobei Bindungen gem\"a\ss{} lokalen Gradienten zur Gibbs-Freie-Energie-Minimierung gebildet und gebrochen werden \citep{MetabolicNash}. Diese Kan\"ale konvergieren zu eng begrenzten lokalen station\"aren Zust\"anden.

\subsection{Mesoebene: Netzwerktopologie und CGT}

Aus der Wechselwirkung lokal gekoppelter Reaktionskan\"ale, die um dasselbe begrenzte Ressourceninventar konkurrieren, entsteht ein kompetitives, frequenzabh\"angiges Spielfeld. Die Regeln dieses Spiels k\"onnen mit CGT beschrieben werden \citep{AzoarcusPNAS}. Auf dieser Ebene bleibt konzentrationsgetriebene Kinetik das Basismodell, muss aber durch bedingte Wahrscheinlichkeiten, Netzwerktopologie und Ljapunow-Stabilit\"at erg\"anzt werden, wenn nicht nur die Anfangsrate, sondern langfristige Persistenz erkl\"art werden soll.

\subsection{Makroebene: MEPP als Stabilit\"atsrandbedingung}

Auf der Makroebene beschr\"ankt das Regime der Nicht-Gleichgewichtsthermodynamik---formalisiert durch MEPP und Englands Modifikation des Zweiten Hauptsatzes---welche lokalen Nash-Gleichgewichte global und langfristig persistieren k\"onnen. Unter der schwachen MEPP-Lesart dieser Arbeit (Abschnitt~\ref{sec:mepp}) muss ein lokal stabiles Netzwerk mit dissipativen Randbedingungen kompatibel sein; ob h\"ohere Entropieproduktion \emph{zwischen} ansonsten lebensf\"ahigen Netzwerken selektiert, ist eine separate empirische Frage. Dieser makroskopische Selektionsclaim \"ubernimmt die in Abschnitt~\ref{sec:cgt} genannten Einschr\"ankungen: Er bleibt eine Arbeitshypothese bis zu den experimentellen Tests aus Abschnitt~\ref{sec:predictions} (P3).

\subsection{\"Uberwindung biologischer Teleologie}

Eine Implikation von FST-II ist die M\"oglichkeit, teleologische Sprache in der Beschreibung biologischer Evolution zu reduzieren \citep{UnbecomingBeing}. Wenn England argumentiert, dass ``der beste Weg, mehr Energie zu dissipieren, einfach darin besteht, mehr Kopien von sich selbst zu machen'', k\"onnen Reproduktion und ``nat\"urliche Selektion'' nicht als inh\"arenter \"Uberlebensinstinkt der DNA, sondern als dissipativ kompatible Prozessbeschreibung gelesen werden.

Leben, mit all seiner makroskopischen Komplexit\"at, seiner molekularen Homochiralit\"at und seinen vernetzten metabolischen Hyperzyklen, kann aus der FST-II-Perspektive als ein \emph{thermodynamisch plausibler} Weg verstanden werden, auf dem Materie in offenen, fern vom Gleichgewicht befindlichen Systemen gro\ss{}e thermische Gradienten dissipiert \citep{Bioenergetics}. Ob dieser Weg \emph{der} wahrscheinlichste oder unvermeidliche ist, bleibt eine offene empirische Frage; der Rahmen beschreibt ihn als Kandidatenattraktor unter MEPP-kompatiblen Bedingungen, beweist aber weder Einzigartigkeit noch Unvermeidlichkeit.

% ===================================================================
\section{Testbare Vorhersagen}
\label{sec:predictions}
% ===================================================================

Tabelle~\ref{tab:prediction-ledger} formuliert die sechs Vorhersagen als
Mess- und Kontrollgates. Sie trennt den vorgeschlagenen spieltheoretischen
Status eines Netzwerks von unabh\"angigen thermodynamischen Observablen und
verhindert, dass die Toy-Rechnungen unten als Validierungsevidenz gelesen
werden.

\begin{table}[htbp]
\centering
\caption{Vorhersage-Mess-Ledger. Jede Vorhersage braucht eine Observable, die
unabh\"angig von dem Payoff-Modell ist, mit dem der spieltheoretische Zustand
definiert wird.}
\label{tab:prediction-ledger}
\footnotesize
\setlength{\tabcolsep}{3pt}
\begin{tabularx}{\textwidth}{@{}>{\raggedright\arraybackslash}p{0.12\textwidth}>{\raggedright\arraybackslash}p{0.21\textwidth}>{\raggedright\arraybackslash}p{0.33\textwidth}>{\raggedright\arraybackslash}X@{}}
\toprule
\textbf{Nr.} & \textbf{Ziel} & \textbf{Unabh\"angige Observablen} & \textbf{Aktuelles Gate} \\
\midrule
P1 & Zykluspersistenz & Kalorimetrisches $\dot{S}_{\mathrm{prod}}$ plus St\"orungserholungszeit f\"ur gematchte Zyklen & Vorgeschlagen; braucht gematchte Zyklen ober- und unterhalb der England-Schwelle \\
P2 & Chiralit\"atsschwelle & $ee(t)$ als Funktion von initialem $ee$, Antriebsst\"arke, autokatalytischen Raten und Kreuzinhibition & Vorgeschlagen; braucht racemische Baseline und Soai-artiges Positivkontrollregime \\
P3 & Netzwerkselektion & Spielstatus aus Sequenzierung oder Payoff-Replay; $\dot{S}_{\mathrm{prod}}$ aus Kalorimetrie & Noch nicht getestet; die bestehende Replikatortabelle ist nur ein Toy-Durchsatzproxy \\
P4 & Kompartiment-Schwelle & Parasitenanteil, Kompartimentgr\"o\ss{}e, Diffusionsl\"ange und Persistenz unter kontrolliertem Fluss & Vorgeschlagen; kompartimentierte und gut durchmischte Kontrollen mit gleicher Feed-Chemie vergleichen \\
P5 & R\"aumlicher Schutz & Invasionsresistenzindex $\eta$ bei identischem Parasiteninput mit und ohne r\"aumliche Struktur & Vorgeschlagen; braucht Gel- oder Oberfl\"achenstruktur-Experimente mit Replikatoren \\
P6 & CGT gegen Kinetik & Anfangswachstumsrate, langfristige Dominanz, Crossover-Zeit und unabh\"angiges $\dot{S}_{\mathrm{prod}}$ & Direktes Falsifikationsgate gegen die Standarderwartung des schnellsten Wachstums \\
\bottomrule
\end{tabularx}
\end{table}

\subsection{P1: Entropieproduktion und Zyklusstabilit\"at}

\begin{proposition}
Unter autokatalytischen Zyklen, die die England-Dissipationsschranke erf\"ullen, weisen diejenigen mit resolvent-stabiler Entropieproduktion (d.\,h.\ st\"orungsresistent in zweiter Ordnung) gr\"o\ss{}ere dynamische Persistenz auf.
\end{proposition}

\textbf{Test:} Konstruiere autokatalytische Zyklen mit variierenden thermodynamischen Effizienzen. Miss die Entropieproduktionsrate (via Kalorimetrie) und die dynamische Stabilit\"at (St\"orungserholungszeit). Die Vorhersage ist keine einfache monotone Korrelation zwischen Dissipationsrate und Stabilit\"at, sondern dass Zyklen oberhalb einer kritischen Dissipationsschwelle einen qualitativen Stabilit\"ats\"ubergang zeigen. Unterhalb der Schwelle (England-Schranke verletzt) zerfallen Zyklen; oberhalb wird Stabilit\"at durch Struktureigenschaften zweiter Ordnung bestimmt (Netzwerktopologie, kreuzkatalytische Kopplung), nicht durch die absolute Dissipationsrate.

\subsection{P2: Chiralit\"atsverst\"arkungs-Dynamik}

\begin{proposition}
Chiralit\"atsverst\"arkung folgt Nash-Gleichgewichts-Dynamik mit einer kritischen Schwelle.
\end{proposition}

\textbf{Test:} In asymmetrisch-autokatalytischen Systemen (z.\,B.\ die Soai-Reaktion \citep{Soai1995}), miss die Verst\"arkungsdynamik des Enantiomeren\"uberschusses (ee) als Funktion des initialen ee und der Antriebsst\"arke (Energiefluss). Der Rahmen sagt vorher: (a) unterhalb einer kritischen Antriebsschwelle bleibt der racemische Zustand stabil (gemischtes Nash-Gleichgewicht); (b) oberhalb dieser Schwelle wird der racemische Zustand zum instabilen Sattelpunkt, und beliebig kleiner initialer ee wird zu Homochiralit\"at verst\"arkt (reines Nash-Gleichgewicht). Die nicht-triviale Vorhersage ist die Existenz und Lage dieser kritischen Antriebsschwelle als Funktion der autokatalytischen Ratenkonstanten und der Kreuzinhibitionsst\"arke.

\subsection{P3: MEPP-Vorhersagen f\"ur Netzwerke}

\begin{proposition}
In Konkurrenzexperimenten zwischen autokatalytischen Netzwerken dominiert unter spezifizierten Anfangsbedingungs- und Stabilit\"atsannahmen ein Netzwerk mit h\"oherer Entropieproduktionsrate.
\end{proposition}

\textbf{Test:} Etabliere Konkurrenz zwischen verschiedenen autokatalytischen Systemen. Miss die Entropieproduktionsraten unabh\"angig, dann f\"uhre die Konkurrenz durch. Die vorhergesagte Tendenz ist, dass das Netzwerk mit h\"oherem $\dot{S}$ das mit niedrigerem $\dot{S}$ verdr\"angt. \textbf{Qualifikation:} Wie das illustrative Beispiel unten zeigt, gilt diese Vorhersage nur langfristig und unter der Bedingung ausreichender anf\"anglicher Kooperatorendichte. In bistabilen Systemen kann das Netzwerk mit h\"oherem $\dot{S}$ bei ung\"unstigen Anfangsbedingungen verlieren (Koordinationsbarriere). Die pr\"azise Vorhersage ist daher: \emph{unter der Bedingung, dass die Koordinationsschwelle \"uberschritten wird}, dominiert das Netzwerk mit h\"oherem $\dot{S}$.

\textbf{Epistemischer Hinweis zur Kausalrichtung.} P3 testet einen \emph{st\"arkeren} Anspruch als die in Abschnitt~\ref{sec:mepp} verteidigte Stabilit\"atsfilter-Lesart: Es sagt voraus, dass Entropieproduktionsrate zwischen lebensf\"ahigen Netzwerken positiv \emph{selektiert}, nicht blo\ss{} die lebensf\"ahige Mannigfaltigkeit abgrenzt. In der reinen Stabilit\"atsfilter-Lesart sollten Netzwerke oberhalb der England-Schranke in Bezug auf $\dot{S}$ selektiv neutral sein; P3 pr\"uft, ob innerhalb der lebensf\"ahigen Mannigfaltigkeit ein verbleibender $\dot{S}$-abh\"angiger Selektionsdruck wirkt. Ein positives Ergebnis w\"urde eine moderate MEPP-als-Selektions-Hypothese st\"utzen; ein Nullergebnis w\"urde die schw\"achere Stabilit\"atsfilter-Interpretation best\"atigen; ein negatives Ergebnis, also robuste Dominanz niedrigerer $\dot{S}$-Netzwerke unter den angegebenen Kontrollen, w\"urde den starken P3-Anspruch falsifizieren, aber nicht die Stabilit\"atsfilter-These auf Framework-Ebene. Alle drei Ergebnisse w\"aren informativ.

\textbf{Illustratives Rechenbeispiel.} Eine Replikator-Dynamik-Simulation mit drei konkurrierenden Netzwerktypen---kooperativ (Azoarcus-artige Kreuzkatalyse, mit frequenzabh\"angiger Fitness $f_C = 1{,}0 + 5{,}0\,x_C^2$), eigenn\"utzig (Selbstkatalyse, konstante Fitness $f_S = 3{,}0$) und parasit\"ar (ausbeuterisch, Fitness $f_P = 0{,}1 + 4{,}0\,x_C$)---enth\"ullt folgende Nash-Gleichgewichtsstruktur. Diese Parameter sind zur Illustration der qualitativen Dynamik gew\"ahlt; sie sind nicht an spezifische experimentelle Daten gefittet:

\begin{table}[ht]
\centering
\caption{Fixpunkte der 3-Spezies-Replikator-Dynamik mit Azoarcus-artigen Fitnessparametern. Die mittlere Fitness $\langle f \rangle$ dient hier als Proxy f\"ur katalytischen Durchsatz; sie ist keine unabh\"angig gemessene thermodynamische Entropieproduktionsrate. In dieser Toy-Parametrisierung erreicht das kooperative Netzwerk den doppelten Durchsatz-Proxy des eigenn\"utzigen Gleichgewichts.}
\label{tab:replicator-fps}
\begin{tabular}{lccc}
\toprule
\textbf{Gleichgewicht} & $\langle f \rangle$ & \textbf{Stabil?} & \textbf{Basin-Gr\"o\ss{}e} \\
\midrule
Kooperativ (ESS) & 6{,}0 & Ja & 41\% \\
Eigenn\"utzig (ESS) & 3{,}0 & Ja & 59\% \\
Parasit\"ar & 0{,}1 & Nein (Sattel) & 0\% \\
\bottomrule
\end{tabular}
\end{table}

Die kooperative ESS hat den doppelten katalytischen Durchsatz-Proxy der eigenn\"utzigen ESS; dies ist qualitativ mit P3 vereinbar, best\"atigt P3 aber nicht. Der Proxy stammt aus dem Payoff-Modell selbst, w\"ahrend P3 eine unabh\"angige thermodynamische Observable wie kalorimetrisches $\dot{S}_{\mathrm{prod}}$ verlangt. Das Beispiel dient daher nur als dynamische Illustration des Koordinationsproblems. Das System zeigt \emph{Bistabilit\"at}: Von einer uniformen Anfangszusammensetzung konvergiert die Dynamik zum eigenn\"utzigen Gleichgewicht, weil der eigenn\"utzige Replikator eine h\"ohere anf\"angliche Wachstumsrate hat ($f_S = 3{,}0$ vs.\ $f_C = 1{,}0 + 5{,}0 \, x_C^2$ f\"ur Kooperatoren). Die kooperative ESS wird nur erreicht, wenn der anf\"angliche Kooperatorenanteil eine kritische Schwelle \"uberschreitet ($x_C \gtrsim 0{,}4$). Diese Koordinationsbarriere erkl\"art, warum Vaidya et al.\ \citep{AzoarcusPNAS} fanden, dass r\"aumliche Struktur (Kompartimentierung, Oberfl\"achenadsorption) f\"ur die Etablierung kooperativer Netzwerke essentiell war---sie konzentriert Kooperatoren lokal \"uber die kritische Schwelle.

\subsection{P4: Kompartimentierungsschwelle}

\begin{proposition}
Die kritische Parasitendichte, ab der Kompartimentierung vorteilhaft wird, kann aus der Hyperzyklus-Spieltheorie berechnet werden.
\end{proposition}

\textbf{Test:} In RNA-Replikator-Experimenten mit variierenden Parasitenanteilen, bestimme die Schwelle, oberhalb derer kompartimentierte Systeme gut durchmischte Systeme \"ubertreffen.

\subsection{P5: Spiralwellen-Schutz}

\begin{proposition}
R\"aumliche Spiralwellen-Organisation sch\"utzt Hyperzyklen gegen parasit\"are Invasion.
\end{proposition}

\textbf{Test:} Vergleiche den Invasionserfolg parasit\"arer RNA-Sequenzen in gut durchmischten vs.\ r\"aumlich strukturierten (gelbasierten) Replikator-Experimenten.

\subsection{P6: Unterscheidung von CGT-Vorhersagen gegen\"uber Standardkinetik}

Ein direkter experimenteller Test, der die spieltheoretische Vorhersage von der Standard-Chemiekinetik unterscheidet, w\"urde zwei konkurrierende autokatalytische Netzwerke unter kontrolliertem Ressourcenfluss umfassen. Die Standardkinetik sagt vorher, dass das Netzwerk mit der h\"oheren anf\"anglichen Wachstumsrate dominiert. Der spieltheoretische Rahmen sagt stattdessen vorher, dass das Netzwerk, das ein resolvent-stabiles Nash-Gleichgewicht erreicht---das m\"oglicherweise eine langsamere Anfangsrate hat, aber robust gegen parasit\"are Invasion ist---langfristig dominieren wird. Dies k\"onnte mit dem Azoarcus-Ribozym-System getestet werden, mit konkurrierenden kooperativen und eigenn\"utzigen RNA-Populationen unter Durchfluss-Bedingungen, wobei sowohl Populationsanteile als auch kalorimetrische Entropieproduktion unabh\"angig gemessen werden. Die Schl\"usselobservable ist, ob das langsamer wachsende kooperative Netzwerk das schneller wachsende eigenn\"utzige Netzwerk nach einer charakteristischen Crossover-Zeit \"uberholt, wie es die Bistabilit\"atsanalyse von P3 vorhersagt.

% ===================================================================
\section{Diskussion}
\label{sec:discussion}
% ===================================================================

\subsection{Beziehung zu Kauffmans autokatalytischen Mengen}

Kauffman \citep{Kauffman1993} schlug vor, dass Leben aus kollektiv autokatalytischen Mengen entstand. FST-II ist kompatibel, f\"ugt aber formale spieltheoretische Struktur hinzu (Nash-Gleichgewichts-Bedingungen), eine thermodynamische Stabilit\"atsrandbedingung (MEPP) und Verbindungen zu Symmetriebrechung und Phasen\"uberg\"angen.

\subsection{Beziehung zu Pross' Dynamic Kinetic Stability}

Pross \citep{Pross2012} f\"uhrte DKS als Persistenzkriterium f\"ur replizierende Systeme ein. Der hier vorgeschlagene Rahmen liefert eine thermodynamische Interpretation: Systeme sind dynamisch stabil, wenn sie das Resolvent-Stabilit\"atskriterium unter dissipativen Randbedingungen erf\"ullen, was in geeigneten chemischen Regimen mit Replikationseffizienz korrelieren kann.

\subsection{Die MEPP-Kontroverse und Geltungsbereich}

Die Verwendung von MEPP als globale Payoff-Funktion erfordert die explizite Anerkennung ihres umstrittenen Status. W\"ahrend MEPP als makroskopisches Organisationsprinzip in der Klimawissenschaft und \"Okosystem-Thermodynamik erfolgreich angewandt wurde \citep{MEPP2006}, sind die theoretischen Grundlagen nicht universell akzeptiert:

\begin{itemize}
\item Grinstein \& Linsker \citep{GrinsteinLinsker2007} zeigten, dass Dewars informationstheoretische Herleitung von MEPP Fehler enth\"alt und MEPP in diskreten stochastischen Systemen versagt.
\item Martyushev \citep{Martyushev2010} identifizierte zwei notwendige Bedingungen f\"ur die Anwendbarkeit von MEPP: (i) lokales thermodynamisches Gleichgewicht und (ii) bilineare Entropieproduktion ($\sigma = \sum_i X_i J_i$). F\"ur die fern vom Gleichgewicht befindlichen autokatalytischen Systeme im Zentrum von FST-II ist Bedingung (ii) in der Regel nicht erf\"ullt.
\item Die in dieser Arbeit entwickelte Resolvent-Perspektive (Abschnitt~\ref{sec:discussion}) adressiert diese Einschr\"ankung: MEPP wird nicht als kausaler Treiber, sondern als \emph{Stabilit\"atsfilter} verwendet, der resolvent-stabile Konfigurationen innerhalb einer degenerierten Mannigfaltigkeit identifiziert. Dieser schw\"achere Anspruch erfordert nicht die volle G\"ultigkeit von MEPP als Selektionsgesetz.
\item Ein aktueller Review von Sawada, Daigaku \& Toma~\citep{Toma2025} best\"atigt, dass MEPP makroskopische evolution\"are Trends---von der Entstehung selbstreplizierender Systeme bis zur sp\"aten gesellschaftlichen Evolution---gut beschreibt, eine rigorose Herleitung aus ersten Prinzipien fern vom Gleichgewicht jedoch unvollst\"andig bleibt. Dies st\"utzt die Positionierung von MEPP als heuristisches Organisationsprinzip in dieser Arbeit.
\end{itemize}

\noindent
Eine begriffliche Kl\"arung ist angebracht. Prigogines Prinzip der \emph{minimalen} Entropieproduktion (mEPP) \citep{Prigogine1947} ist im linearen Regime nahe dem Gleichgewicht rigoros bewiesen, wo Onsagersche Reziprozit\"atsrelationen gelten und die Entropieproduktion bilinear in Kr\"aften und Fl\"ussen ist: station\"are Zust\"ande minimieren $\sigma$. Das Prinzip der \emph{maximalen} Entropieproduktion (MEPP), verbunden mit Ziegler \citep{Ziegler1963} und zusammengefasst von Martyushev \& Seleznev \citep{MEPP2006}, postuliert hingegen, dass Systeme fern vom Gleichgewicht Zust\"ande w\"ahlen, die $\sigma$ maximieren---eine Heuristik ohne rigorose Herleitung f\"ur allgemeine nichtlineare Systeme. Das scheinbare Paradoxon (min vs.\ max) l\"ost sich durch Regime-Trennung auf: mEPP gilt lokal f\"ur Subsysteme nahe ihren individuellen station\"aren Zust\"anden, w\"ahrend MEPP als Selektionsprinzip zwischen konkurrierenden Konfigurationen auf der Makroskala fern vom Gleichgewicht wirkt. Im FST-Kontext ist diese Unterscheidung wesentlich: die autokatalytischen Netzwerke aus Abschnitt~\ref{sec:hypercycles} operieren fern vom Gleichgewicht, wo mEPP nicht gilt, und MEPP tritt nur als der oben beschriebene Stabilit\"atsfilter ein.

\subsection{Was ist wirklich neu?}

\begin{enumerate}
    \item \textbf{Formale Synthese:} Vereinigung der chemischen Gibbs-Gleichgewicht--Nash-Gleichgewicht-Analogie \citep{MetabolicNash} mit der England-Dissipationsschranke \citep{England2013} und der Eigen-Schuster-Replikator-Dynamik \citep{Eigen1979} in einen einzigen koh\"arenten Rahmen---keine dieser Verbindungen wurde zuvor zusammengef\"uhrt
    \item \textbf{MEPP-Neuinterpretation:} Die Resolvent-Stabilit\"atsfilter-Lesart von MEPP, die die kausale Ambiguit\"at in der Entropieproduktions-Literatur aufl\"ost
    \item \textbf{Spieltheoretische Chiralit\"at:} Neuinterpretation von Franks Symmetriebrechung \citep{Frank1953} als \"Ubergang vom gemischten zum reinen Nash-Gleichgewicht, was ein Vokabular f\"ur quantitative Vorhersagen (P2) liefert
    \item \textbf{Skalen\"ubergreifende Architektur:} Integration in den FST-Rahmen, der Teilchenphysik (FST-I), Chemie (FST-II) und Biologie (FST-III) durch eine wiederkehrende Pattern-A-Stabilit\"atslogik verbindet---diese skalen\"ubergreifende strukturelle Korrespondenz ist der prim\"are originelle Beitrag
\end{enumerate}

Diese Liste ist bewusst begrenzt. FST-II beansprucht nicht, Autokatalyse, Hyperzyklen, Viedma-Ripening, Frank-artige Chiralit\"atsverst\"arkung, Replikator-Dynamik oder die Standardgesetze chemischer Raten neu zu entdecken. Diese Mechanismen stammen aus der etablierten Literatur. Der vorgeschlagene Beitrag liegt in der dar\"ubergelegten Ordnungsschicht: einer spieltheoretischen Darstellung von Netzwerkkonkurrenz, einer Stabilit\"atsfilter-Lesart dissipativer Randbedingungen und experimentell trennbaren Vorhersagen (insbesondere P3 und P6), die scheitern k\"onnen, falls Standardkinetik allein die langfristig dominierenden Netzwerke erkl\"art.

Die mathematische Verbindung zwischen Eigen-Schuster-Stabilit\"at und Nash-Stabilit\"at ist im folgenden eingeschr\"ankten Sinne exakt: \emph{an Fixpunktl\"osungen der Replikatorgleichung} ist jedes Nash-Gleichgewicht der Payoff-Matrix $\mathbf{A}$ ein station\"arer Punkt, und jeder asymptotisch stabile station\"are Punkt ist ein Nash-Gleichgewicht (Theorem~\ref{thm:replicator-nash}). Diese Korrespondenz ist gut etabliert \citep{HofbauerSigmund2003}.

\subsection{Stabilit\"at vs.\ Maximierung: Die Resolventenperspektive}

Der FST-Rahmen beschreibt \emph{Stabilit\"atsprinzipien}, keine Maximierungsprinzipien. Die Dynamik f\"uhrender Ordnung ist entartet: viele Konfigurationen pr\"abiotischer Netzwerke sind in erster Ordnung gleicherma\ss{}en lebensf\"ahig. Die Selektion erfolgt in zweiter Ordnung \"uber resolventenartige Kopplung zwischen entarteten und komplement\"aren Unterr\"aumen. MEPP wirkt hier als Stabilit\"atsfilter, der einen resolventenstabilen Fixpunkt innerhalb der entarteten Mannigfaltigkeit identifizieren soll. Diese strukturelle Logik---als strukturelle Analogie in der eichtheoretischen Reformulierung der Riemannschen Vermutung entwickelt \citep{GeigerRH2026b}---hilft, die MEPP-Kontroverse (das Prinzip selektiert f\"ur Stabilit\"at, nicht f\"ur maximale Dissipation), das Problem der kausalen Richtung von Englands Ungleichung (Dissipationsschranken markieren die lebensf\"ahige Mannigfaltigkeit; Resolventenstabilit\"at selektiert innerhalb dieser) und das Chiralit\"atsproblem (Homochiralit\"at kann als resolventenstabiler Fixpunkt eines in f\"uhrender Ordnung entarteten Gleichgewichts gelesen werden) gemeinsam zu adressieren. \"Uber alle drei FST-Arbeiten hinweg gilt die gleiche Selektionslogik zweiter Ordnung: Nash-Gleichgewichte sind keine Optima, sondern resolventenstabile Fixpunkte.

Im vorliegenden Kontext veranschaulicht das Racemat diese Architektur: in f\"uhrender Ordnung haben L- und D-Enantiomere identische thermodynamische Payoffs (die entartete Mannigfaltigkeit). MEPP selektiert den homochiralen Zustand nicht, weil er mehr Entropie produziert---in erster Ordnung dissipiert ein Racemat vergleichbar. Vielmehr destabilisiert die Kopplung zweiter Ordnung durch autokatalytische Verst\"arkung und Kreuzinhibition die racemische Konfiguration und selektiert den homochiralen Zustand als einen resolventenstabilen Fixpunkt im Frank-Modell.

\paragraph{Operationalisierung der Resolvent-Stabilit\"at.}
Das Resolvent-Stabilit\"atskriterium l\"asst eine konkrete chemische
Interpretation zu.  Man betrachte ein linearisiertes Reaktions-Diffusions-System
am station\"aren Konzentrationsprofil~$\mathbf{c}^{*}$.  Die Jacobi-Matrix
$A = \partial \mathbf{f}/\partial \mathbf{c}\big|_{\mathbf{c}^{*}}$
kodiert die lokale Kopplung aller Spezies.  F\"ur offene Systeme ohne exakte
Erhaltungsgesetze erfordert Resolvent-Stabilit\"at, dass jeder Eigenwert von~$A$
strikt negativen Realteil besitzt.  F\"ur konservative Reaktionsnetzwerke muss
dieses Kriterium zuerst auf die st\"ochiometrische Kompatibilit\"atsklasse (oder
\"aquivalent auf den Tangentialraum nach Quotientierung erhaltener Gr\"o\ss{}en)
eingeschr\"ankt werden; neutrale Massenerhaltungsmoden sind keine
Instabilit\"atsmoden.  In beiden F\"allen wird die Stabilit\"atsgrenze durch die
Spektralabszisse $\alpha(A)=\max \operatorname{Re}\lambda(A)$ auf dem relevanten
Unterraum kontrolliert, nicht durch den Spektralradius allein.  Der station\"are
Zustand ist asymptotisch stabil, wenn diese eingeschr\"ankte Spektralabszisse
negativ ist und Trajektorien nach hinreichend kleinen zul\"assigen St\"orungen zu
$\mathbf{c}^{*}$ zur\"uckkehren.
Diese Bedingung ist experimentell \"uber mindestens zwei Wege zug\"anglich:
(i)~\emph{Relaxationszeitmessungen} nach einer kontrollierten Perturbation
(Verd\"unnungspuls, Temperatursprung, pH-Shift), wobei die langsamste
exponentielle Abklingrate direkt den am wenigsten negativen Eigenwert liefert;
und (ii)~\emph{systematische St\"orungsscans}, die abbilden, wie sich die
eingeschr\"ankte Spektralabszisse~$\alpha(A)$ \"uber Umgebungsparameter der Null
n\"ahert und so die Grenze der resolvent-stabilen Region im Parameterraum
kartieren. Der Spektralradius~$\rho(A)$ bleibt als Steifheits- oder
Oszillationsskala n\"utzlich, ist f\"ur sich allein aber kein Zertifikat einer
Stabilit\"atsmarge.

In diesem Bild markiert Englands Dissipationsschranke~\citep{England2013} die
\emph{lebensf\"ahige Mannigfaltigkeit}: Nur Netzwerke, deren Entropieproduktion
das thermodynamische Minimum \"ubersteigt, k\"onnen \"uberhaupt persistieren
(Lemma~\ref{lem:england}).  Die Selektion zweiter Ordnung wirkt dann
\emph{innerhalb} dieser Mannigfaltigkeit.  Parasitenresistenz, r\"aumliche
Struktur (Spiralwellen in Hyperzyklen) und Ressourcenfluss-Topologie bestimmen,
\emph{welches} der lebensf\"ahigen Netzwerke resolvent-stabil ist und daher die
Langzeitdynamik dominiert.
Diese Zwei-Stufen-Architektur -- Lebensf\"ahigkeitsfilter erster Ordnung plus
Resolventenselektion zweiter Ordnung -- verbindet die formale Stabilit\"atstheorie
des FST-RH-Rahmens~\citep{GeigerRH2026b} mit Gr\"o\ss{}en, die direkt im
chemischen Labor messbar sind.

\textbf{Die behauptete strukturelle Korrespondenz zwischen Gibbs-Minimierung und Nash-Gleichgewichten bedarf jedoch der Qualifizierung.} Gibbs-Freie-Energie-Minimierung (Abschnitt 3.2) ist ein konvexes Optimierungsproblem mit einem einzigen globalen Minimum im thermodynamischen Limes. Nash-Gleichgewichte m\"ussen hingegen nicht einzigartig sein, k\"onnen nicht-konvex sein und Sattelpunkte einschlie\ss{}en. Die Behauptung, dass ``ein formales Nash-Gleichgewicht genau dann erreicht wird, wenn kein Reaktionspfad seine Gibbs-Energie weiter minimieren kann'', verwechselt lokale Optimalit\"atsbedingungen (Stationarit\"at erster Ordnung) mit Nash-Gleichgewicht im spieltheoretischen Sinne. Was tats\"achlich gezeigt wird, ist, dass \emph{die station\"aren Punkte des eingeschr\"ankten Gibbs-Minimierungsproblems die notwendigen Bedingungen f\"ur ein Nash-Gleichgewicht erf\"ullen}---dies ist eine formale Analogie, kein vollst\"andiger Isomorphismus. Die globale Eindeutigkeit des thermodynamischen Gleichgewichts \"ubertr\"agt sich nicht auf die spieltheoretische Situation, wo mehrere Nash-Gleichgewichte koexistieren k\"onnen.

\subsection{Computational Outlook: Machine Learning Interatomic Potentials}

Ein vielversprechender Weg zur quantitativen Verifikation der hier vorgeschlagenen chemischen Spielmatrizen liegt in Machine-Learning-basierten interatomaren Potenzialen (MLIPs). Ein aktueller Review von Anton \& Stirnemann~\citep{AntonStirnemann2026} zeigt, dass neuronale Netzwerk-Potenziale, die auf Ab-initio-Daten trainiert wurden, autokatalytische Reaktionsnetzwerke, \"Ubergangszust\"ande und Reaktionskinetiken mit nahezu quantenmechanischer Genauigkeit bei klassischen Rechenkosten simulieren k\"onnen. F\"ur das FST-II-Framework k\"onnten MLIPs die Absch\"atzung der Ratenkonstanten automatisieren, die in die Payoff-Matrizen der Abschnitte~\ref{sec:cgt}--\ref{sec:hypercycles} eingehen, mittels Active-Learning-Workflows, die Potenzialenergieoberfl\"achen entlang relevanter Reaktionskoordinaten iterativ verfeinern. Insbesondere k\"onnten Freie-Energie-Profile f\"ur konkurrierende autokatalytische Zyklen---derzeit nur durch aufwendige Dichtefunktionaltheorie-Molekulardynamik zug\"anglich---auf Mikrosekunden-Zeitskalen berechnet werden, was eine direkte numerische Auswertung der Nash-Gleichgewichts-Vorhersagen P1--P3 erm\"oglicht. Ein solcher Ansatz w\"urde auch ein systematisches Screening von Umgebungsparametern (pH, Temperatur, Mineraloberfl\"achen) erlauben, um die resolvent-stabilen Regionen der pr\"abiotischen Spiellandschaft zu kartieren. Obwohl MLIPs die volle Komplexit\"at offener, fern vom Gleichgewicht befindlicher Reaktionsnetzwerke noch nicht erfassen, legt der rasante methodische Fortschritt, der in~\citep{AntonStirnemann2026} zusammengefasst wird, nahe, dass eine rechnerische Validierung der spieltheoretischen Vorhersagen innerhalb der n\"achsten Jahre m\"oglich werden k\"onnte.

\subsection{Vorgeschlagener nicht-zirkul\"arer experimenteller Test}
\label{sec:noncircular-test}

Das in Abschnitt~\ref{sec:cgt} identifizierte Zirkularit\"atsrisiko (MEPP definiert gleichzeitig die Payoff-Funktion und das vorhergesagte Ergebnis) kann empirisch getestet werden, indem Entropieproduktion und spieltheoretischer Status \"uber \emph{unabh\"angige} Observablen gemessen werden. Wir schlagen folgendes Protokoll vor:

\begin{enumerate}
    \item \textbf{System:} Zwei konkurrierende autokatalytische RNA-Netzwerke---z.B.\ Azoarcus-Ribozym-Varianten mit unterschiedlicher Substratpr\"aferenz \citep{AzoarcusPMC}---werden in einem gemeinsamen Kompartiment mit begrenzten Ressourcen platziert.
    \item \textbf{Entropieproduktion (Observable 1):} Die Gesamtdissipationsrate wird \emph{separat} mittels isothermer Titrationskalorimetrie (ITC) gemessen, die die W\"armefreisetzung pro Zeiteinheit quantifiziert, oder \"uber eine geschlossene Redox-Bilanz / Photonenfluss-Buchf\"uhrung bei photogetriebenen Systemen. Dies liefert $\dot{S}_{\mathrm{prod}}(t)$ ohne Bezug auf die Replikator-Identit\"at.
    \item \textbf{Spielstatus (Observable 2):} Die Populationszusammensetzung---welches Netzwerk dominiert, ob Koexistenz oder Oszillation vorliegt---wird \emph{unabh\"angig} durch RNA-Sequenzierung oder Fluoreszenzmarkierung netzwerkspezifischer Motive bestimmt.
    \item \textbf{Korrelationstest:} Die FST-II-Vorhersage ist, dass das resolvent-stabile Nash-Gleichgewicht (die langfristig gewinnende Netzwerkkonfiguration) mit h\"oherer zeitgemittelter Entropieproduktion $\langle\dot{S}_{\mathrm{prod}}\rangle$ korreliert. Da Spielstatus und Entropieproduktion aus verschiedenen Messkan\"alen stammen (Sequenzierung vs.\ Kalorimetrie), w\"urde eine positive Korrelation eine nicht-zirkul\"are empirische St\"utzung unter den angegebenen Kontrollen liefern.
\end{enumerate}

\noindent
Dieses Design adressiert direkt Einschr\"ankung (L3) und liefert den von Vorhersage P3 geforderten unabh\"angigen Test.

\subsection{\"Uber Azoarcus hinaus: Kandidatensysteme f\"ur FST-II-Tests}
\label{sec:beyond-azoarcus}

Die empirische Basis von FST-II beruht prim\"ar auf dem Azoarcus-Ribozym-System. Um Allgemeing\"ultigkeit zu etablieren, eignen sich mindestens drei weitere molekulare Plattformen zum Testen der Vorhersagen dieses Rahmens:

\begin{itemize}
    \item \textbf{Autokatalytische Peptidnetzwerke.} Die Ashkenasy-Gruppe hat selbstreplizierende Peptidnetzwerke demonstriert, die kooperative und kompetitive Dynamiken analog zu den vier Azoarcus-Spielen zeigen \citep{Ashkenasy2004}. Peptidsysteme bieten den Vorteil einstellbarer Interaktionsst\"arken \"uber Sequenzmutation, was eine systematische Variation der Payoff-Matrix-Eintr\"age erm\"oglicht.
    \item \textbf{Lipidvesikel--Autokatalysator-Systeme.} Arbeiten aus dem Szostak-Labor zu kompartimentierten Selbstreplikatoren \citep{Szostak2001} liefern ein nat\"urliches Testfeld f\"ur Vorhersage P4 (Kompartimentierungsschwelle): Vesikel-eingekapselte autokatalytische Zyklen k\"onnen mit Kontrollexperimenten in freier L\"osung verglichen werden, um zu testen, ob Kompartimentierung das Nash-Gleichgewicht in Richtung kooperativer Ergebnisse verschiebt, wie FST-II vorhersagt.
    \item \textbf{Alternative RNA-Ribozyme.} Die R3C-Ligase \citep{Rogers2001} und Hammerhead-Ribozyme \citep{Hammerhead1986} bieten strukturell verschiedene selbstspaltende/ligierende Motive. Kreuzkatalytische Paare aus diesen Familien k\"onnten verwendet werden, um Spielmatrizen unabh\"angig von der Azoarcus-Topologie zu konstruieren und zu testen, ob die FST-II-Stabilit\"atsvorhersagen f\"ur verschiedene Netzwerkarchitekturen gelten.
\end{itemize}

\noindent
F\"ur jedes System ist die Schl\"usselanforderung, dass (i) Nash-Gleichgewichts-artige Dynamiken (Dominanz, Koexistenz, Bistabilit\"at) beobachtbar sind und (ii) die Entropieproduktion unabh\"angig gemessen werden kann---dasselbe Dual-Observable-Design wie in Abschnitt~\ref{sec:noncircular-test} vorgeschlagen.

\subsection{Nicht-ideale Korrekturen}
\label{sec:nonideal-corrections}

Die Gibbs-Freie-Energie-Formulierung in Abschnitt~\ref{sec:cgt} nimmt ideale Mischung ($\gamma_i = 1$) an. Drei Klassen nicht-idealer Korrekturen werden f\"ur realistische pr\"abiotische Szenarien relevant:

\begin{enumerate}
    \item \textbf{Aktivit\"atskoeffizienten.} Bei hohen Konzentrationen oder f\"ur geladene Spezies (z.B.\ RNA-Oligomere mit mehreren Phosphatgruppen) weichen die Aktivit\"atskoeffizienten $\gamma_i$ erheblich von eins ab. Die Gibbs-Energie lautet dann $G = \sum_i n_i (G_i^0 + RT\ln(\gamma_i n_i / n_{\mathrm{total}}))$, und die Nash-Gleichgewichtszusammensetzungen verschieben sich entsprechend. Erweiterte Debye--H\"uckel- oder Pitzer-Modelle liefern systematische Korrekturen f\"ur ionische L\"osungen.
    \item \textbf{R\"aumliche Heterogenit\"at.} Konzentrationsgradienten (z.B.\ nahe Mineraloberfl\"achen oder innerhalb von Protozell-Kompartimenten) modifizieren die effektive Payoff-Matrix: lokale Konzentrationen unterscheiden sich von Durchschnittswerten, sodass das Spiel, das ein Replikator erf\"ahrt, von seiner r\"aumlichen Position abh\"angt. Dies verwandelt die gut durchmischte Replikatorgleichung in ein Reaktions-Diffusions-Spiel, in dem r\"aumliche ESS (Abschnitt~\ref{sec:hypercycles}) und Kompartimentierungseffekte (P4) interagieren.
    \item \textbf{Auswirkung auf die Stabilit\"atsanalyse.} Beide Korrekturen gehen in die lineare Stabilit\"atsanalyse zweiter Ordnung ein: wenn $\gamma_i \neq 1$, erh\"alt die Jacobi-Matrix der Replikatordynamik zus\"atzliche Terme proportional zu $\partial\gamma_i/\partial c_j$, die Eigenwerte verschieben und den resolvent-stabilen Fixpunkt modifizieren k\"onnen. Die qualitative spieltheoretische Struktur (Existenz von Nash-Gleichgewichten) bleibt erhalten, aber die quantitativen Vorhersagen (welches Gleichgewicht selektiert wird, kritische Schwellenwerte f\"ur P1--P4) erfordern Neuberechnung mit den korrigierten Payoff-Funktionen.
\end{enumerate}

\noindent
Diese nicht-idealen Effekte verkn\"upfen sich mit der Aktivit\"atskoeffizienten-Bemerkung in Abschnitt~\ref{sec:cgt} und mit Einschr\"ankung (L5). Eine systematische Behandlung wird auf zuk\"unftige Arbeit verschoben, ist aber essenziell f\"ur den quantitativen Vergleich mit Experimenten bei realistischen pr\"abiotischen Konzentrationen.

\subsection{Minimales Rechenbeispiel: Drei-Spezies-Autokatalyse}\label{sec:3species}

Zur numerischen Illustration des Zusammenhangs zwischen Resolventenstabilit\"at und Entropieproduktion analysieren wir einen minimalen autokatalytischen Drei-Spezies-Zyklus auf dem Konzentrationssimplex $\Delta^2$.
Die Reaktionen lauten $\mathrm{A}+\mathrm{B}\xrightarrow{k_1}2\mathrm{B}$ ($k_1=1{,}0$), $\mathrm{B}+\mathrm{C}\xrightarrow{k_2}2\mathrm{C}$ ($k_2=0{,}8$), $\mathrm{C}\xrightarrow{k_3}\mathrm{A}$ ($k_3=0{,}5$).
Dabei ist $\mathrm{C}\to\mathrm{A}$ kein autokatalytischer Schritt, sondern ein Rezyklierungs- bzw.\ Abbauschritt erster Ordnung; er schlie\ss{}t den Zyklus, indem er das Substrat~A regeneriert, und verhindert damit den trivialen absorbierenden Zustand $x_C=1$. Die Asymmetrie zwischen den autokatalytischen Schritten ($\mathrm{A}+\mathrm{B}\to 2\mathrm{B}$, $\mathrm{B}+\mathrm{C}\to 2\mathrm{C}$) und diesem Abbauschritt ist beabsichtigt: Sie modelliert die biologisch verbreitete Situation, dass in einem zyklischen Netzwerk nicht alle Schritte katalytisch sind.
Das System besitzt einen eindeutigen inneren Fixpunkt $\mathbf{x}^*=(0{,}167;\,0{,}625;\,0{,}208)$, der die Nash-Gleichgewichtsbedingung des assoziierten Replikatorspiels erf\"ullt, mit Gibbs-Freier-Energie $G(\mathbf{x}^*)=-0{,}711$ (in dimensionslosen Einheiten).

Da die Dynamik die Gesamtkonzentration auf dem Simplex erh\"alt, besitzt die volle Jacobi-Matrix einen neutralen Massenerhaltungsmodus. Beschr\"ankt auf den Tangentialraum des Konzentrationssimplex haben die nicht-neutralen Eigenwerte negative Realteile (n\"aherungsweise $-0{,}3125 \pm 0{,}2997\,i$) und best\"atigen asymptotische Stabilit\"at innerhalb des physikalisch relevanten Simplex. Der berichtete Spektralradius $\rho(J)=0{,}433$ ist daher am besten als lokale Steifheits-/Oszillationsskala zu lesen, nicht als w\"ortlicher Abstand zur Stabilit\"atsgrenze.
Die Schnakenberg-Entropieproduktionsrate am Fixpunkt betr\"agt $\dot{S}=1{,}439$ \citep{Schnakenberg1976}.
Um die Beziehung zwischen Resolventenstabilit\"at und Dissipation zu untersuchen, wurde $k_1$ \"uber 20 Werte variiert und $\rho(J)$ sowie $\dot{S}$ am jeweiligen Fixpunkt berechnet.
Die Pearson-Korrelation betr\"agt $r=0{,}801$ und zeigt eine positive Assoziation zwischen Spektralradius und Entropieproduktionsrate. \textbf{Hinweis:} Da sowohl $\rho(J)$ als auch $\dot{S}$ deterministische Funktionen des einzigen variierten Parameters~$k_1$ sind, ist diese Korrelation eine mathematische Konsequenz der parametrischen Abh\"angigkeit und kein unabh\"angiger statistischer Befund. Der $p$-Wert eines Pearson-Tests ist in diesem Ein-Parameter-Kontext epistemisch nicht aussagekr\"aftig. Die Korrelation illustriert eine parameterabh\"angige Steifheits--Dissipations-Kopplung, stellt aber keine Evidenz f\"ur ihre Universalit\"at \"uber strukturell verschiedene Netzwerke hinweg dar (siehe Hinweis unten).

\begin{table}[htbp]
\centering
\caption{Zusammenfassung der Drei-Spezies-Autokatalyse-Berechnung (Skript: \texttt{autocatalysis\_3species.py}).}
\label{tab:3species}
\small
\begin{tabular}{ll}
\toprule
\textbf{Parameter} & \textbf{Wert} \\
\midrule
$x^*_A,\; x^*_B,\; x^*_C$ & $0{,}167;\; 0{,}625;\; 0{,}208$ \\
$\rho(J)$ & $0{,}433$ \\
$\dot{S}$ & $1{,}439$ \\
Pearson $r(\rho,\dot{S})$ & $0{,}801$ (Ein-Parameter-Scan; siehe Text) \\
\bottomrule
\end{tabular}
\end{table}

Diese Machbarkeitsrechnung liefert keinen allgemeinen Beweis, macht aber eine lokale Kovarianz sichtbar, die f\"ur die FST-II-Architektur relevant ist: schnellerer autokatalytischer Umsatz erh\"oht zugleich die Entropieproduktionsrate und ver\"andert die lokale Jacobi-Skala. In diesem Toy-Modell w\"achst $\rho(J)$ mit $\dot{S}$, aber $\rho(J)$ ist ein Steifheits-/Oszillationsproxy und kein Zertifikat f\"ur eine Stabilit\"atsmarge.
Die starke parametrische Korrelation ($r>0{,}8$) ist mit einem m\"oglichen Steifheits--Dissipations-Zielkonflikt vereinbar und rechtfertigt eine systematische Untersuchung gr\"o\ss{}erer und strukturell vielf\"altiger Netzwerke, bei der nicht nur eine einzelne Ratenkonstante, sondern die Netzwerktopologie variiert wird, wie in Abschnitt~\ref{sec:predictions} skizziert.

\textbf{Hinweis zur kausalen Interpretation.} Die berichtete Korrelation entsteht durch die Variation eines einzigen Parameters~($k_1$); sowohl $\rho(J)$ als auch $\dot{S}$ sind Funktionen von~$k_1$, sodass eine hohe Korrelation aus rein mathematischen Gr\"unden zu erwarten ist und f\"ur sich genommen keine physikalische Beziehung zwischen Stabilit\"at und Dissipation belegt. Ein echter Stabilit\"atsmargentest m\"usste den maximalen Realteil der Jacobi-Matrix im Tangentialraum oder den Bifurkationsabstand unter kontrollierten St\"orungen verfolgen, nicht $\rho(J)$ allein. Ein strengerer Test m\"usste die \emph{Netzwerktopologie} (Anzahl der Spezies, Konnektivit\"at, relative Gr\"o\ss{}en aller Ratenkonstanten) variieren und pr\"ufen, ob die Steifheits--Dissipations-Relation auch \"uber strukturell verschiedene Netzwerke hinweg bestehen bleibt. Der vorliegende Ein-Parameter-Scan sollte daher als Illustration eines m\"oglichen Mechanismus gelesen werden, nicht als Evidenz f\"ur seine Universalit\"at.

Die beobachtete Kovarianz ist mit dem FST-II-Rahmen vereinbar: Englands Dissipationsschranke markiert die lebensf\"ahige Mannigfaltigkeit, und eine an MEPP orientierte Stabilit\"atsfilterung wird innerhalb der resolvent-stabilen Region hypothetisch angesetzt. Die positive Korrelation zwischen $\rho(J)$ und $\dot{S}$ ist mit einem stabilen Arbeitspunkt in diesem Toy-Modell vereinbar; sie belegt f\"ur sich allein aber weder einen kausalen Druck zu h\"oherer Entropieproduktion noch eine Bewegung in Richtung einer Stabilit\"atsgrenze. Der st\"arkere Selektionsclaim erfordert Topologie-Tests \"uber verschiedene Netzwerke hinweg, wie sie in Vorhersage~P1 und im nicht-zirkul\"aren Protokoll von Abschnitt~\ref{sec:noncircular-test} spezifiziert sind.

\subsection{Status der Behauptungen}

Tabelle~\ref{tab:status-claims} klassifiziert die zentralen Behauptungen dieses Artikels nach epistemischem Status.

\begin{table}[htbp]
\centering
\caption{Status der Behauptungen. \textsc{Etabliert}: bewiesene oder weithin akzeptierte Ergebnisse, die hier zitiert werden; \textsc{Konjektur}: neue, aber testbare Behauptungen; \textsc{Spekulativ}: Rahmenvorschl\"age ohne direkten empirischen Test.}
\label{tab:status-claims}
\small
\setlength{\tabcolsep}{4pt}
\begin{tabularx}{\textwidth}{@{}>{\raggedright\arraybackslash}p{0.29\textwidth}>{\raggedright\arraybackslash}p{0.20\textwidth}>{\raggedright\arraybackslash}X@{}}
\toprule
\textbf{Behauptung} & \textbf{Status} & \textbf{Evidenz / Referenz} \\
\midrule
Replikator--Nash-Korrespondenz & \textsc{Etabliert} & Hofbauer \& Sigmund (1998, 2003); Theorem~\ref{thm:replicator-nash} \\
England-Dissipationsschranke & \textsc{Etabliert} & England (2013); Crooks (1999); Lemma~\ref{lem:england} \\
MEPP als heuristisches Organisationsprinzip & \textsc{Etabliert} & Martyushev \& Seleznev (2006); umstritten fern vom Gleichgewicht \\
Azoarcus-Ribozyme als spieltheoretische Systeme & \textsc{Etabliert} & Vaidya et al.\ (2012); Sievers \& von Kiedrowski (1994) \\
Frank-Modell f\"ur Chiralit\"atsverst\"arkung & \textsc{Etabliert} & Frank (1953); Viedma-Ripening-Experimente \\
Gibbs-Minimierung--Nash-Gleichgewicht-Analogie & \textsc{Konjektur} & Abschn.~\ref{sec:cgt}; formale Analogie, kein Isomorphismus \\
Resolventenstabilit\"at als Selektionsmechanismus & \textsc{Konjektur} & Remark~\ref{rem:autocatalysis-resolvent}; strukturelle Analogie, keine dom\"anenspezifische Formalisierung \\
Homochiralit\"at als Reine-Strategie-Nash-\"Ubergang & \textsc{Konjektur} & Abschn.~\ref{sec:chirality}; Neuinterpretation von Frank (1953) \\
MEPP-Netzwerkwettbewerb (P3) & \textsc{Konjektur} & Abschn.~\ref{sec:predictions}; illustrative Simulation, kein Experiment \\
Spiralwellen als verteilte ESS & \textsc{Konjektur} & Abschn.~\ref{sec:hypercycles}; qualitativ, kein quantitativer DESS-Beweis \\
Kompartimentierungs-Schwellenwert (P4) & \textsc{Konjektur} & Testbar; noch keine experimentellen Daten \\
Hierarchische Mikro/Meso/Makro-Kaskade & \textsc{Spekulativ} & Abschn.~\ref{sec:synthesis}; konzeptuelle Architektur \\
\bottomrule
\end{tabularx}
\end{table}

% ===================================================================
\section{Schlussfolgerung}
\label{sec:conclusion}
% ===================================================================

Der FST-II-Rahmen liefert ein strukturiertes Interpretationsmodell der pr\"abiotischen chemischen Evolution. Durch die mathematische Verbindung der Ljapunow-Metrik differentieller Replikatorgleichungen mit thermodynamischer Entropieproduktion \citep{ReplicatorEntropy} legt er nahe, dass die Entstehung des Lebens als m\"ogliche Folge von Stabilit\"atsselektion unter dissipativen Randbedingungen in offenen materiellen Systemen untersucht werden kann.

Pr\"abiologische Reaktionen sind keine isolierten stochastischen Kollisionen allein, sondern vernetzte Prozesse, die als thermodynamische Mehrspieler-Spiele repr\"asentiert werden k\"onnen. Der zentrale ``Payoff'' in diesen molekularen Spielen ist ein Kandidatenma\ss{} f\"ur stabilit\"atsselektierte Dissipation. Ph\"anomene wie Homochiralit\"at und autokatalytische Hyperzyklen k\"onnen als Nash-Gleichgewichts-artige Strukturen innerhalb dieser energetischen Wettbewerbe interpretiert werden.

Mehrere wichtige Einschr\"ankungen sind zu beachten:
\begin{enumerate}[label=(L\arabic*)]
    \item \textbf{MEPP-Status:} MEPP wird als heuristischer Stabilit\"atsfilter verwendet, nicht als bewiesenes kausales Prinzip. Die Bedingungen, unter denen MEPP auf fern vom Gleichgewicht befindliche autokatalytische Systeme anwendbar ist, m\"ussen noch rigoros etabliert werden (Abschnitt~\ref{sec:discussion}).
    \item \textbf{Empirischer Geltungsbereich:} Die empirische Basis beruht prim\"ar auf dem Azoarcus-Ribozym-System (16 Genotypen, spezifische In-vitro-Bedingungen). Verallgemeinerung auf andere Molek\"ulklassen (Peptide, Lipide) und Umweltbedingungen erfordert unabh\"angige experimentelle Validierung.
    \item \textbf{Zirkularit\"atsrisiko:} Die Identifikation von MEPP als Payoff-Funktion, die Nash-Gleichgewichte definiert, erzeugt eine Definitionsschleife (Abschnitt~\ref{sec:cgt}). Unabh\"angige Messung von Entropieproduktion und spieltheoretischem Status (P3) ist erforderlich, um diese Zirkularit\"at aufzubrechen.
    \item \textbf{Resolvententerminologie:} Das Resolventenstabilit\"ats-Vokabular wird als strukturelle Analogie zum FST-RH-Rahmen verwendet; eine dom\"anenspezifische mathematische Formulierung f\"ur chemische Systeme bleibt zu entwickeln.
    \item \textbf{Idealit\"atsannahme:} Die Gibbs-Freie-Energie-Formulierung nimmt ideale Mischung an (Abschnitt~\ref{sec:cgt}); reale pr\"abiotische Systeme erfordern Korrekturen der Aktivit\"atskoeffizienten.
\end{enumerate}
Als Rahmenartikel pr\"asentiert FST-II keine neuen experimentellen Daten und keine neuen mathematischen Beweise; sein Beitrag ist die Synthese etablierter Ergebnisse zu einer koh\"arenten Architektur, die testbare Vorhersagen erzeugt. Die sechs testbaren Vorhersagen (P1--P6) liefern ein konkretes experimentelles Programm, um diesen Rahmen von einer blo\ss{}en Umformulierung bekannter Ergebnisse zu unterscheiden. Der Wert des Rahmens wird letztlich daran zu messen sein, ob diese Vorhersagen durch k\"unftige Experimente best\"atigt, pr\"azisiert oder widerlegt werden.

Trotz dieser Einschr\"ankungen er\"offnet FST-II vorsichtige, bislang ungetestete M\"oglichkeiten, erwartbare exobiologische Systemarchitekturen auf anderen Planeten einzugrenzen und als konzeptioneller Leitfaden f\"ur das rationale Design k\"unstlicher dissipativer Netzwerke in der modernen synthetischen Systemchemie zu dienen.

% ===================================================================
% Literaturverzeichnis
% ===================================================================

\bibliographystyle{plainnat}
\begin{thebibliography}{99}

\bibitem[Geiger(2026a)]{FSTI}
Geiger, L. (2026).
\newblock Functional Stability Theory I: A Game-Theoretic Framework for the Thermodynamic Stability of Fundamental Parameters.
\newblock \textit{Zenodo preprint}, Version 1.1, 2026.
\newblock DOI: 10.5281/zenodo.20130955.

\bibitem[Frank(1953)]{Frank1953}
Frank, F.~C. (1953).
\newblock On spontaneous asymmetric synthesis.
\newblock \textit{Biochimica et Biophysica Acta}, 11, 459--463.
\newblock DOI: 10.1016/0006-3002(53)90082-1.

\bibitem[Hofbauer \& Sigmund(1998)]{HofbauerSigmund1998}
Hofbauer, J. \& Sigmund, K. (1998).
\newblock \textit{Evolutionary Games and Population Dynamics}.
\newblock Cambridge University Press.

\bibitem[Hofbauer \& Sigmund(2003)]{HofbauerSigmund2003}
Hofbauer, J. \& Sigmund, K. (2003).
\newblock Evolutionary game dynamics.
\newblock \textit{Bulletin of the American Mathematical Society}, 40(4), 479--519.
\newblock DOI: 10.1090/S0273-0979-03-00988-1.

\bibitem[Yeates et al.(2016)]{AzoarcusPNAS}
Yeates, J.~A.~M., Hilbe, C., Zwick, M., Nowak, M.~A. \& Lehman, N. (2016).
\newblock Dynamics of prebiotic RNA reproduction illuminated by chemical game theory.
\newblock \textit{PNAS}, 113(18), 5030--5035.
\newblock DOI: 10.1073/pnas.1525273113.

\bibitem[Vaidya et al.(2012)]{AzoarcusPMC}
Vaidya, N., Manapat, M.~L., Chen, I.~A., Xulvi-Brunet, R., Hayden, E.~J. \& Lehman, N. (2012).
\newblock Spontaneous network formation among cooperative RNA replicators.
\newblock \textit{Nature}, 491, 72--77.
\newblock DOI: 10.1038/nature11549.

\bibitem[Dyakin \& Uversky(2022)]{Biochirality}
Dyakin, V.~V. \& Uversky, V.~N. (2022).
\newblock Arrow of time, entropy, and protein folding: Holistic view on biochirality.
\newblock \textit{International Journal of Molecular Sciences}, 23(7), 3687.
\newblock DOI: 10.3390/ijms23073687.

\bibitem[Nicolis \& Prigogine(1977)]{Bioenergetics}
Nicolis, G. \& Prigogine, I. (1977).
\newblock \textit{Self-Organization in Non-Equilibrium Systems: From Dissipative Structures to Order Through Fluctuations}.
\newblock Wiley-Interscience.

\bibitem[Velegol et al.(2018)]{CGT2018}
Velegol, D., Suhey, P., Connolly, J., Morrissey, N. \& Cook, L. (2018).
\newblock Chemical game theory.
\newblock \textit{Industrial \& Engineering Chemistry Research}, 57(41), 13593--13607.
\newblock DOI: 10.1021/acs.iecr.8b03835.

\bibitem[Nowak(2006)]{CGTMacmillan}
Nowak, M.~A. (2006).
\newblock \textit{Evolutionary Dynamics: Exploring the Equations of Life}.
\newblock Harvard University Press.

\bibitem[Eigen \& Schuster(1979)]{Eigen1979}
Eigen, M. \& Schuster, P. (1979).
\newblock \textit{The Hypercycle}.
\newblock Springer.

\bibitem[England(2013)]{England2013}
England, J.~L. (2013).
\newblock Statistical physics of self-replication.
\newblock \textit{The Journal of Chemical Physics}, 139, 121923.
\newblock DOI: 10.1063/1.4818538.

\bibitem[Martyushev(2013)]{EntropyMEPP}
Martyushev, L.~M. (2013).
\newblock Entropy and entropy production: Old misconceptions and new breakthroughs.
\newblock \textit{Entropy}, 15(4), 1152--1170.
\newblock DOI: 10.3390/e15041152.

% Entropianism.org reference removed (non-scholarly source).

\bibitem[Viedma(2005)]{HomochiralityThermo}
Viedma, C. (2005).
\newblock Chiral symmetry breaking during crystallization: Complete chiral purity induced by nonlinear autocatalysis and recycling.
\newblock \textit{Physical Review Letters}, 94(6), 065504.
\newblock DOI: 10.1103/PhysRevLett.94.065504.

\bibitem[Boerlijst \& Hogeweg(1991)]{HypercycleInfinite}
Boerlijst, M.~C. \& Hogeweg, P. (1991).
\newblock Spiral wave structure in pre-biotic evolution: hypercycles stable against parasites.
\newblock \textit{Physica D: Nonlinear Phenomena}, 48(1), 17--28.
\newblock DOI: 10.1016/0167-2789(91)90049-F.

\bibitem[Kauffman(1993)]{Kauffman1993}
Kauffman, S.~A. (1993).
\newblock \textit{The Origins of Order: Self-Organization and Selection in Evolution}.
\newblock Oxford University Press.
\newblock DOI: 10.1093/oso/9780195079517.001.0001.

\bibitem[Grinstein \& Linsker(2007)]{GrinsteinLinsker2007}
Grinstein, G. \& Linsker, R. (2007).
\newblock Comments on a derivation and application of the `maximum entropy production' principle.
\newblock \textit{Journal of Physics A: Mathematical and Theoretical}, 40(31), 9717--9720.
\newblock DOI: 10.1088/1751-8113/40/31/N01.

\bibitem[Martyushev(2010)]{Martyushev2010}
Martyushev, L.~M. (2010).
\newblock The maximum entropy production principle: two basic questions.
\newblock \textit{Philosophical Transactions of the Royal Society B}, 365(1545), 1333--1334.
\newblock DOI: 10.1098/rstb.2010.0110.

\bibitem[Martyushev \& Seleznev(2006)]{MEPP2006}
Martyushev, L.~M. \& Seleznev, V.~D. (2006).
\newblock Maximum entropy production principle in physics, chemistry and biology.
\newblock \textit{Physics Reports}, 426(1), 1--45.
\newblock DOI: 10.1016/j.physrep.2005.12.001.

\bibitem[Schuster et al.(2008)]{MetabolicNash}
Schuster, S., Kreft, J.-U., Schroeter, A. \& Pfeiffer, T. (2008).
\newblock Use of game-theoretical methods in biochemistry and biophysics.
\newblock \textit{Journal of Biological Physics}, 34(1--2), 1--17.
\newblock DOI: 10.1007/s10867-008-9101-4.

\bibitem[Miller(1953)]{MillerUrey}
Miller, S.~L. (1953).
\newblock A production of amino acids under possible primitive Earth conditions.
\newblock \textit{Science}, 117(3046), 528--529.
\newblock DOI: 10.1126/science.117.3046.528.

% OvershootLoop reference (Odum & Pinkerton, 1955) removed: no longer cited in text after attribution correction.

\bibitem[Prigogine(1947)]{Prigogine1947}
Prigogine, I. (1947).
\newblock \textit{\'Etude thermodynamique des ph\'enom\`enes irr\'eversibles}.
\newblock Desoer, Li\`ege.

\bibitem[Pross(2012)]{Pross2012}
Pross, A. (2012).
\newblock \textit{What is Life? How Chemistry Becomes Biology}.
\newblock Oxford University Press.

% Key renamed from SchmitzReview2011 (v1) to PlassonReview2007 (v2) to match actual reference.
\bibitem[Plasson et al.(2007)]{PlassonReview2007}
Plasson, R., Kondepudi, D.~K., Bersini, H., Commeyras, A. \& Asakura, K. (2007).
\newblock Emergence of homochirality in far-from-equilibrium systems: Mechanisms and role in prebiotic chemistry.
\newblock \textit{Chirality}, 19(8), 589--600.
\newblock DOI: 10.1002/chir.20440.

\bibitem[Hutson \& Vickers(1992)]{ReplicatorDiffusion}
Hutson, V.~C.~L. \& Vickers, G.~T. (1992).
\newblock Travelling waves and dominance of ESS's.
\newblock \textit{Journal of Mathematical Biology}, 30(5), 457--471.
\newblock DOI: 10.1007/BF00160531.

\bibitem[Harper(2009)]{ReplicatorEntropy}
Harper, M. (2009).
\newblock The replicator equation as an inference dynamic.
\newblock arXiv:0911.1763.

\bibitem[Boerlijst \& Hogeweg(1995)]{SpiralWaves}
Boerlijst, M.~C. \& Hogeweg, P. (1995).
\newblock Spatial gradients enhance persistence of hypercycles.
\newblock \textit{Physica D: Nonlinear Phenomena}, 88(1), 29--39.
\newblock DOI: 10.1016/0167-2789(95)00178-7.

% StrategicPersistence merged into Pross2012 (same book, duplicate removed).

\bibitem[Perunov et al.(2016)]{ThermodynamicSelection}
Perunov, N., Marsland, R.~A. \& England, J.~L. (2016).
\newblock Statistical physics of adaptation.
\newblock \textit{Physical Review X}, 6(2), 021036.
\newblock DOI: 10.1103/PhysRevX.6.021036.

\bibitem[Soai et al.(1995)]{Soai1995}
Soai, K., Shibata, T., Morioka, H. \& Choji, K. (1995).
\newblock Asymmetric autocatalysis and amplification of enantiomeric excess of a chiral molecule.
\newblock \textit{Nature}, 378, 767--768.
\newblock DOI: 10.1038/378767a0.

\bibitem[Swenson(1989)]{ThermodynamicsLife}
Swenson, R. (1989).
\newblock Emergent attractors and the law of maximum entropy production: Foundations to a theory of general evolution.
\newblock \textit{Systems Research}, 6(3), 187--197.
\newblock DOI: 10.1002/sres.3850060302.

\bibitem[De Bari et al.(2023)]{ThermodynamicsOrganisms}
De Bari, B., Dixon, J., Kondepudi, D.~K. \& Vaidya, A. (2023).
\newblock Thermodynamics, organisms and behaviour.
\newblock \textit{Philosophical Transactions of the Royal Society A}, 381(2252), 20220278.
\newblock DOI: 10.1098/rsta.2022.0278.

\bibitem[Kondepudi \& Nelson(1985)]{TimeReversalSymmetry}
Kondepudi, D.~K. \& Nelson, G.~W. (1985).
\newblock Weak neutral currents and the origin of biomolecular chirality.
\newblock \textit{Nature}, 314, 438--441.
\newblock DOI: 10.1038/314438a0.

\bibitem[Carroll(2016)]{UnbecomingBeing}
Carroll, S.~M. (2016).
\newblock \textit{The Big Picture: On the Origins of Life, Meaning, and the Universe Itself}.
\newblock Dutton/Penguin.

\bibitem[Gould(1989)]{Gould1989}
Gould, S.~J. (1989).
\newblock \textit{Wonderful Life: The Burgess Shale and the Nature of History}.
\newblock W.~W. Norton.

\bibitem[Crooks(1999)]{Crooks1999}
Crooks, G.~E. (1999).
\newblock Entropy production fluctuation theorem and the nonequilibrium work relation for free energy differences.
\newblock \textit{Physical Review E}, 60(3), 2721--2726.
\newblock DOI: 10.1103/PhysRevE.60.2721.

\bibitem[Geiger(2026c)]{GeigerRH2026c}
Geiger, L. (2026).
\newblock From Landscape to Proof: The Even-Dominance Route to the Riemann Hypothesis.
\newblock \textit{Zenodo preprint}, Version 2.1, 2026.
\newblock DOI: 10.5281/zenodo.19546593.

\bibitem[Geiger(2026b)]{GeigerRH2026b}
Geiger, L. (2026).
\newblock A Proof of the Riemann Hypothesis via Even Dominance of the Weil Quadratic Form.
\newblock \textit{Zenodo preprint}, Version 2.0, 2026.
\newblock DOI: 10.5281/zenodo.19536076.

\bibitem[Geiger(2026d)]{FST-Unified2026}
Geiger, L. (2026).
\newblock FST Spectrum Duality: The Renormalized Free Energy Principle as Physical Instantiation of Functional Stability Theory.
\newblock \textit{Zenodo preprint}, Version 1.7, 2026.
\newblock DOI: 10.5281/zenodo.19162705.

\bibitem[Ziegler(1963)]{Ziegler1963}
Ziegler, H. (1963).
\newblock Some extremum principles in irreversible thermodynamics with application to continuum mechanics.
\newblock In: Sneddon, I.~N. \& Hill, R. (Eds.), \textit{Progress in Solid Mechanics}, Vol.~4, pp.~91--193. North-Holland, Amsterdam.

\bibitem[Weibull(1995)]{Weibull1995}
Weibull, J.~W. (1995).
\newblock \textit{Evolutionary Game Theory}.
\newblock MIT Press.

\bibitem[Sawada et al.(2025)]{Toma2025}
Sawada, Y., Daigaku, Y. \& Toma, K. (2025).
\newblock Maximum entropy production principle of thermodynamics for the birth and evolution of life.
\newblock \textit{Entropy}, 27(4), 449.
\newblock DOI: 10.3390/e27040449.

\bibitem[Anton \& Stirnemann(2026)]{AntonStirnemann2026}
Anton, O. \& Stirnemann, G. (2026).
\newblock Computational studies of prebiotic chemistry at the age of machine learning: From recent breakthroughs to future revolutions.
\newblock \textit{ChemSystemsChem}, 8(1), e00057.
\newblock DOI: 10.1002/syst.202500057.

\bibitem[Ashkenasy et al.(2004)]{Ashkenasy2004}
Ashkenasy, G., Jagasia, R., Yadav, M. \& Ghadiri, M.~R. (2004).
\newblock Design of a directed molecular network.
\newblock \textit{Proceedings of the National Academy of Sciences}, 101(30), 10872--10877.
\newblock DOI: 10.1073/pnas.0402674101.

\bibitem[Szostak et al.(2001)]{Szostak2001}
Szostak, J.~W., Bartel, D.~P. \& Luisi, P.~L. (2001).
\newblock Synthesizing life.
\newblock \textit{Nature}, 409, 387--390.
\newblock DOI: 10.1038/35053176.

\bibitem[Rogers \& Joyce(2001)]{Rogers2001}
Rogers, J. \& Joyce, G.~F. (2001).
\newblock The effect of cytidine on the structure and function of an RNA ligase ribozyme.
\newblock \textit{RNA}, 7(3), 395--404.
\newblock DOI: 10.1017/S135583820100228X.

\bibitem[Prody et al.(1986)]{Hammerhead1986}
Prody, G.~A., Bakos, J.~T., Buzayan, J.~M., Schneider, I.~R. \& Bruening, G. (1986).
\newblock Autolytic processing of dimeric plant virus satellite RNA.
\newblock \textit{Science}, 231(4745), 1577--1580.
\newblock DOI: 10.1126/science.231.4745.1577.

\bibitem[Gilbert(1986)]{Gilbert1986}
Gilbert, W. (1986).
\newblock Origin of life: The RNA world.
\newblock \textit{Nature}, 319, 618.
\newblock DOI: 10.1038/319618a0.

\bibitem[Sievers \& von Kiedrowski(1994)]{SieversKiedrowski1994}
Sievers, D. \& von Kiedrowski, G. (1994).
\newblock Self-replication of complementary nucleotide-based oligomers.
\newblock \textit{Nature}, 369, 221--224.
\newblock DOI: 10.1038/369221a0.

\bibitem[Blackmond(2010)]{Blackmond2010}
Blackmond, D.~G. (2010).
\newblock The origin of biological homochirality.
\newblock \textit{Cold Spring Harbor Perspectives in Biology}, 2(5), a002147.
\newblock DOI: 10.1101/cshperspect.a002147.

\bibitem[Schnakenberg(1976)]{Schnakenberg1976}
Schnakenberg, J. (1976).
\newblock Network theory of microscopic and macroscopic behavior of master equation systems.
\newblock \textit{Reviews of Modern Physics}, 48(4), 571--585.
\newblock DOI: 10.1103/RevModPhys.48.571.

\end{thebibliography}

\section*{KI-Offenlegung}
Diese Arbeit wurde von Claude (Anthropic) bei der Literatursynthese, mathematischen Verifikation und Texterstellung unterst\"utzt. Alle wissenschaftlichen Behauptungen, originalen Argumente und theoretischen Beitr\"age liegen in der alleinigen Verantwortung des menschlichen Autors. Kein KI-System ist als Koautor gelistet.

\end{document}
